% ----------------------------------------------------------------------
%  3.2.2. Зонная теория
% ----------------------------------------------------------------------
\tikzset{
  bzfront/.style={atomLine, line width=0.7pt, line join=round},
  bzback/.style={atomLine!55, line width=0.4pt, dashed},
  bzpt/.style={circle, fill=fmRed, inner sep=0pt, minimum size=3pt},
}
\subsection{Зонная теория}\label{ssec:bands}

В реальном кристалле ионы расположены периодически, и электроны движутся
не в ровном «ящике», а в периодически меняющемся поле решётки. Зонная
теория, в отличие от теории свободных электронов, учитывает это поле.
Итог окажется простым и важным: разрешённые энергии электрона
группируются в \emph{зоны}, разделённые \emph{запрещёнными промежутками}.
Именно расположение и заполнение зон решает, будет ли кристалл металлом,
полупроводником или изолятором.

\paragraph{Периодический потенциал}
Поле отдельного иона в металле действует на меньших расстояниях, чем
кулоновское поле такого же изолированного иона: коллективизированные
валентные электроны \emph{экранируют} заряд иона. Истинный потенциал
$U(\vb r)$, в котором движется электрон,~— периодическое повторение
экранированных потенциалов ионов (рис.~\ref{fig:periodic-pot},\,\textit{а}).
Движение электрона описывает уравнение Шрёдингера
\begin{equation}\label{eq:schr-crystal}
\left[-\frac{\hbar^2}{2m}\nabla^2+U(\vb r)\right]\psi=E\psi,
\qquad U(\vb r+\vb a)=U(\vb r),
\end{equation}
где $\vb a=n_1\vb a_1+n_2\vb a_2+n_3\vb a_3$~— любой вектор трансляции
решётки ($n_i$~— целые).

\begin{figure}[ht]
\centering
\begin{tikzpicture}[x=1cm, y=1cm]
  % а — экранированные потенциалы ионов
  \begin{scope}
    \draw[-{Stealth[length=2mm]}, fmGray] (-0.3,0) -- (6.6,0) node[right, font=\small]{$x$};
    \draw[fmBlue, thick, smooth, samples=160, domain=-0.2:6.2] plot (\x,
      {-0.35/(abs(\x-0.5)+0.12)-0.35/(abs(\x-2)+0.12)-0.35/(abs(\x-3.5)+0.12)-0.35/(abs(\x-5)+0.12)});
    \foreach \x in {0.5,2,3.5,5}{\node[atom=2.6mm] at (\x,0.45) {};}
    \node[font=\small, fmBlue] at (6.2,-1.1) {$U(x)$};
    \node[font=\small] at (3,-3.8) {\textit{а}};
  \end{scope}
  % б — потенциал Кронига–Пенни
  \begin{scope}[shift={(7.6,0)}]
    \draw[-{Stealth[length=2mm]}, fmGray] (-0.3,-2.6) -- (6.4,-2.6) node[right, font=\small]{$x$};
    \draw[fmBlue, thick] (-0.2,-2.6) -- (0,-2.6) -- (0,-0.6) -- (0.4,-0.6) -- (0.4,-2.6) -- (1.6,-2.6) -- (1.6,-0.6)
      -- (2,-0.6) -- (2,-2.6) -- (3.2,-2.6) -- (3.2,-0.6) -- (3.6,-0.6) -- (3.6,-2.6) -- (4.8,-2.6) -- (4.8,-0.6)
      -- (5.2,-0.6) -- (5.2,-2.6) -- (5.9,-2.6);
    \draw[dimline] (0.4,-3.0) -- node[below, font=\small]{$a$} (1.6,-3.0);
    \draw[dimline] (1.6,-3.0) -- node[below, font=\small]{$b$} (2.0,-3.0);
    \draw[dimline] (2.3,-2.6) -- node[right, font=\small]{$U_0$} (2.3,-0.6);
    \draw[fmRed, dashed] (-0.2,-1.7) -- (5.9,-1.7) node[right, font=\small]{$E$};
    \node[font=\small] at (3,-3.8) {\textit{б}};
  \end{scope}
\end{tikzpicture}
\caption{Потенциальная энергия электрона вдоль ряда ионов:
\textit{а}~— реальный периодический потенциал (ямы у ионов);
\textit{б}~— упрощённая модель Кронига--Пенни: прямоугольные ямы шириной
$a$, разделённые барьерами высотой $U_0>E$ и шириной $b$}
\label{fig:periodic-pot}
\end{figure}

\paragraph{Теорема Блоха}
Раз потенциал периодичен, то и вероятность найти электрон должна быть
периодичной: $|\psi(\vb r+\vb a)|^2=|\psi(\vb r)|^2$~— поведение
электрона в одной ячейке ничем не отличается от поведения в любой
другой. Отсюда следует (это общее следствие трансляционной симметрии,
п.~1.2) \term{теорема Блоха}: волновая функция электрона в кристалле
имеет вид
\begin{keyeq}\label{eq:bloch}
\psi_{\vb k}(\vb r)=u_{\vb k}(\vb r)\,e^{i\vb k\vb r},\qquad
u_{\vb k}(\vb r+\vb a)=u_{\vb k}(\vb r).
\end{keyeq}
Это плоская бегущая волна, как у свободного электрона, но
«промодулированная» периодической функцией $u_{\vb k}$ с периодом
решётки. Такие функции называют \term{функциями Блоха}.

\begin{glavnoe}
Электрон Блоха с одинаковой вероятностью находится в эквивалентных
точках \emph{каждой} элементарной ячейки: он не привязан к одному иону,
а принадлежит кристаллу в целом. Блоховская волна проходит через
\emph{идеальную} решётку, не рассеиваясь. Это и есть ответ на загадку,
с которой начался раздел: идеальный кристалл «прозрачен» для
электронов. Сопротивление металла создают только отклонения от
идеальной периодичности~— тепловые колебания атомов, примеси, дефекты
(п.~3.5).
\end{glavnoe}

\paragraph{Модель Кронига--Пенни}
Найти энергии электрона в произвольном периодическом поле трудно. Но
главные черты видны уже на простейшей одномерной модели Кронига и
Пенни (рис.~\ref{fig:periodic-pot},\,\textit{б}): ямы шириной $a$ с
нулевой потенциальной энергией чередуются с барьерами высотой $U_0>E$ и
шириной $b$.
\begin{shagi}
\item Внутри ямы решение~— сумма волн $e^{\pm i\alpha x}$, где
$\alpha=\sqrt{2mE}/\hbar$; внутри барьера~— сумма экспонент $e^{\pm\beta
x}$, $\beta=\sqrt{2m(U_0-E)}/\hbar$.
\item Функция и её производная должны быть непрерывны на границах ям и
барьеров, а решение в целом~— иметь вид Блоха. Это четыре линейных
однородных уравнения для четырёх коэффициентов.
\item Ненулевое решение существует, если определитель системы равен
нулю. Это даёт связь энергии $E$ (через $\alpha$ и $\beta$) с волновым
числом $k$.
\end{shagi}

\begin{podrobnee}[условие разрешимости]
Определитель обращается в нуль при
\[
\frac{\beta^2-\alpha^2}{2\alpha\beta}\,\operatorname{sh}(\beta b)\sin(\alpha a)+
\operatorname{ch}(\beta b)\cos(\alpha a)=\cos\bigl[k(a+b)\bigr].
\]
Упростим, перейдя к пределу узких и высоких барьеров: $b\to0$,
$U_0\to\infty$ при постоянной «мощности» барьера $U_0b$. Тогда
$\operatorname{ch}(\beta b)\to1$, $\operatorname{sh}(\beta b)\to\beta b$,
$\beta^2\gg\alpha^2$, и условие принимает вид \eqref{eq:kp}.
\end{podrobnee}

В пределе тонких высоких барьеров получается уравнение Кронига--Пенни:
\begin{keyeq}\label{eq:kp}
P\,\frac{\sin\alpha a}{\alpha a}+\cos\alpha a=\cos ka,\qquad
P=\frac{mabU_0}{\hbar^2}.
\end{keyeq}
Параметр $P$ характеризует «непрозрачность» барьеров: при $P=0$
барьеров нет и электроны свободны, при $P\to\infty$ ямы полностью
изолированы друг от друга (как атомы, удалённые на большое расстояние).

\paragraph{Разрешённые и запрещённые энергии}
Правая часть \eqref{eq:kp}, $\cos ka$, лежит между $-1$ и $+1$. Значит,
решения существуют только при тех $\alpha a$, при которых левая часть,
$f(\alpha a)$, тоже не выходит за эти пределы. Построим $f(\alpha a)$
(рис.~\ref{fig:kp-f}): там, где $|f|>1$, уравнение Шрёдингера решений не
имеет~— это \term{запрещённые зоны}; отрезки $\Delta_1,\Delta_2,\ldots$,
где $|f|\le1$,~— \term{разрешённые зоны}. Так как $\alpha a$ однозначно
связано с энергией ($E=\hbar^2\alpha^2/2m$), каждая разрешённая полоса
$\alpha a$ даёт полосу разрешённых энергий.

% сгенерировано tools/gen_bands.py — не править вручную
\def\KPloA{2.2505}\def\KPhiA{3.1416}
\def\KPloB{4.7124}\def\KPhiB{6.2832}
\def\KPloC{7.4154}\def\KPhiC{9.4248}
\def\KPloD{10.2841}\def\KPhiD{12.5664}

\begin{figure}[ht]
\centering
\begin{tikzpicture}
\begin{axis}[width=0.92\textwidth, height=5.8cm, xmin=0, xmax=12.9, ymin=-2, ymax=4.2,
  xlabel={$\alpha a$}, ylabel={$f(\alpha a)$},
  xtick={3.14159,6.28319,9.42478,12.56637}, xticklabels={$\pi$,$2\pi$,$3\pi$,$4\pi$},
  ytick={-1,0,1}, axis lines=left, clip=true]
  \fill[fmGreen!14] (axis cs:\KPloA,-2) rectangle (axis cs:\KPhiA,4.2);
  \fill[fmGreen!14] (axis cs:\KPloB,-2) rectangle (axis cs:\KPhiB,4.2);
  \fill[fmGreen!14] (axis cs:\KPloC,-2) rectangle (axis cs:\KPhiC,4.2);
  \fill[fmGreen!14] (axis cs:\KPloD,-2) rectangle (axis cs:\KPhiD,4.2);
  \draw[fmGray, dashed] (axis cs:0,1) -- (axis cs:12.9,1);
  \draw[fmGray, dashed] (axis cs:0,-1) -- (axis cs:12.9,-1);
  \draw[fmGray] (axis cs:0,0) -- (axis cs:12.9,0);
  \addplot[fmBlue, thick] table {figures/kp-f.dat};
  \draw[fmGreen!60!black, line width=2.4pt] (axis cs:\KPloA,0) -- (axis cs:\KPhiA,0);
  \draw[fmGreen!60!black, line width=2.4pt] (axis cs:\KPloB,0) -- (axis cs:\KPhiB,0);
  \draw[fmGreen!60!black, line width=2.4pt] (axis cs:\KPloC,0) -- (axis cs:\KPhiC,0);
  \draw[fmGreen!60!black, line width=2.4pt] (axis cs:\KPloD,0) -- (axis cs:\KPhiD,0);
  \node[font=\small, fmGreen!50!black] at (axis cs:2.7,-1.6) {$\Delta_1$};
  \node[font=\small, fmGreen!50!black] at (axis cs:5.5,-1.6) {$\Delta_2$};
  \node[font=\small, fmGreen!50!black] at (axis cs:8.4,-1.6) {$\Delta_3$};
  \node[font=\small, fmGreen!50!black] at (axis cs:11.4,-1.6) {$\Delta_4$};
\end{axis}
\end{tikzpicture}
\caption{Графическое решение уравнения Кронига--Пенни при $P=3\pi/2$.
Решения есть только там, где $|f(\alpha a)|\le1$ (зелёные полосы
$\Delta_1\ldots\Delta_4$~— разрешённые зоны); между ними~— запрещённые
зоны. С ростом энергии разрешённые зоны расширяются, а запрещённые
сужаются}
\label{fig:kp-f}
\end{figure}

\paragraph{Закон дисперсии $E(k)$}
Внутри первой разрешённой зоны $\cos ka$ меняется от $+1$ до $-1$, т.\,е.
$k$ пробегает значения от 0 до $\pm\pi/a$. Внутри второй зоны $\cos ka$
снова проходит от $-1$ до $+1$~— этому соответствуют $\pi/a<|k|<2\pi/a$,
и так далее. В результате кривая $E(k)$ состоит из кусков, близких к
параболе свободного электрона $E=\hbar^2k^2/2m$, но \emph{разорванных}
при
\begin{keyeq}\label{eq:gap-k}
k=\frac{n\pi}{a},\qquad n=\pm1,\pm2,\ldots
\end{keyeq}
(рис.~\ref{fig:kp-ek}). Вблизи разрывов кривая отклоняется от параболы
и подходит к границе \emph{горизонтально}: групповая скорость электрона
$v=\hbar^{-1}\d E/\d k$ обращается в нуль~— решётка «тормозит» волну.

\begin{figure}[ht]
\centering
\begin{tikzpicture}
\begin{axis}[width=0.86\textwidth, height=7.4cm, xmin=-4.2, xmax=4.2, ymin=0, ymax=165,
  xlabel={$k$}, ylabel={$E$, ед. $\hbar^2/2ma^2$},
  xtick={-4,-3,-2,-1,0,1,2,3,4},
  xticklabels={$-\frac{4\pi}a$,$-\frac{3\pi}a$,$-\frac{2\pi}a$,$-\frac\pi a$,0,$\frac\pi a$,$\frac{2\pi}a$,$\frac{3\pi}a$,$\frac{4\pi}a$},
  ytick=\empty, axis lines=left, clip=false]
  \fill[fmRed!8] (axis cs:-4.2,9.87) rectangle (axis cs:4.2,22.2);
  \fill[fmRed!8] (axis cs:-4.2,39.48) rectangle (axis cs:4.2,54.99);
  \fill[fmRed!8] (axis cs:-4.2,88.83) rectangle (axis cs:4.2,105.8);
  \addplot[fmGray, dashed, domain=-4.1:4.1, samples=100] {(pi*x)^2};
  \addplot[fmBlue, very thick] table {figures/kp-band1.dat};
  \addplot[fmBlue, very thick] table {figures/kp-band1m.dat};
  \addplot[fmBlue, very thick] table {figures/kp-band2.dat};
  \addplot[fmBlue, very thick] table {figures/kp-band2m.dat};
  \addplot[fmBlue, very thick] table {figures/kp-band3.dat};
  \addplot[fmBlue, very thick] table {figures/kp-band3m.dat};
  \addplot[fmBlue, very thick] table {figures/kp-band4.dat};
  \addplot[fmBlue, very thick] table {figures/kp-band4m.dat};
  \pgfplotsinvokeforeach{-3,-2,-1,1,2,3}{\draw[fmGray!60, dotted] (axis cs:#1,0) -- (axis cs:#1,165);}
  \node[font=\scriptsize, fmRed, anchor=west, align=left] at (axis cs:4.25,47) {запрещённые\\зоны};
  \node[font=\scriptsize, fmGray] at (axis cs:0,140) {пунктир: $E=\hbar^2k^2/2m$};
\end{axis}
\end{tikzpicture}
\caption{Энергия электрона в модели Кронига--Пенни ($P=3\pi/2$, расчёт).
Синие кривые~— разрешённые зоны, красные полосы~— запрещённые; пунктир~—
парабола свободного электрона. Разрывы возникают при $k=n\pi/a$}
\label{fig:kp-ek}
\end{figure}

\paragraph{Почему возникают разрывы: брэгговское отражение}
Электрон на уровне Ферми ($E\sim5$~эВ) имеет длину волны
$\lambda=2\pi\hbar/\sqrt{2mE}\approx0{,}55$~нм~— того же порядка, что и
межатомные расстояния (0,2--0,3~нм). Кристалл служит для электронной
волны дифракционной решёткой. Подставим $k=2\pi/\lambda$ в условие
\eqref{eq:gap-k}: $n\lambda=2a$. Это \emph{условие Вульфа--Брэгга}
$n\lambda=2d\sin\theta$ для волны, падающей перпендикулярно отражающим
плоскостям ($\theta=90^\circ$, $d=a$).

Такая волна отражается назад, а падающая и отражённая волны дают
\emph{стоячую} волну. Их можно сложить двумя способами
(рис.~\ref{fig:standing}):
\[
\psi_+\propto\cos\frac{\pi x}a,\qquad \psi_-\propto\sin\frac{\pi x}a .
\]
В состоянии $\psi_+$ электронная плотность $|\psi_+|^2$ сосредоточена на
ионах, где потенциальная энергия наиболее низка; в состоянии $\psi_-$~—
между ионами. Энергии двух состояний различаются на $\Delta E$~— это и
есть ширина запрещённой зоны. Электроны с волновыми числами,
удовлетворяющими условию Брэгга, не могут бежать сквозь решётку.

\begin{figure}[ht]
\centering
\begin{tikzpicture}
\begin{axis}[width=0.86\textwidth, height=5.4cm, xmin=-0.3, xmax=3.3, ymin=-1.9, ymax=2.35,
  axis lines=none, clip=false]
  \addplot[fmGray, thin, domain=-0.3:3.3, samples=200]
    {max(-1.85,-0.12/(abs(x)+0.06)-0.12/(abs(x-1)+0.06)-0.12/(abs(x-2)+0.06)-0.12/(abs(x-3)+0.06))};
  \addplot[fmBlue, thick, domain=-0.3:3.3, samples=200] {2*cos(deg(pi*x))^2};
  \addplot[fmRed, thick, domain=-0.3:3.3, samples=200] {2*sin(deg(pi*x))^2};
  \addplot[black, dashed, domain=-0.3:3.3] {1};
  \pgfplotsinvokeforeach{0,1,2,3}{\node[atom=2.6mm] at (axis cs:#1,-0.2) {};}
  \node[font=\small, fmBlue, anchor=west] at (axis cs:3.32,1.95) {$|\psi_+|^2$};
  \node[font=\small, fmRed, anchor=west] at (axis cs:3.32,0.3) {$|\psi_-|^2$};
  \node[font=\small, anchor=west] at (axis cs:3.32,1.0) {бегущая волна};
  \node[font=\small, fmGray, anchor=west] at (axis cs:3.32,-1.4) {$U(x)$};
  \draw[dimline] (axis cs:0,2.25) -- node[above, font=\small]{$a$} (axis cs:1,2.25);
\end{axis}
\end{tikzpicture}
\caption{Две стоячие волны на границе зоны ($k=\pi/a$): у $\psi_+$
электронная плотность скапливается на ионах (энергия ниже), у
$\psi_-$~— между ионами (энергия выше). У бегущей волны плотность
одинакова везде (пунктир)}
\label{fig:standing}
\end{figure}

\begin{smysl}[Второй взгляд: от атомных уровней к зонам]
К зонам можно прийти и с другой стороны~— от изолированных атомов. Пока
$N$ атомов далеко друг от друга, каждый их уровень повторяется $N$ раз.
При сближении волновые функции соседей перекрываются, и уровень
расщепляется на $N$ близких подуровней~— образуется зона
(рис.~\ref{fig:levels-bands}). Чем сильнее перекрытие, тем шире зона:
внутренние оболочки почти не перекрываются и дают узкие зоны, внешние
(валентные) уровни расплываются в широкие зоны, которые могут даже
перекрываться. Это предел $P\to\infty$ модели Кронига--Пенни, когда ямы
почти изолированы.
\end{smysl}

\begin{figure}[ht]
\centering
\begin{tikzpicture}[x=2.2cm, y=0.55cm]
  \draw[-{Stealth[length=2mm]}] (0.6,-9.2) -- (4.3,-9.2) node[right, font=\small]{$r$};
  \draw[-{Stealth[length=2mm]}] (0.6,-9.2) -- (0.6,1.2) node[above, font=\small]{$E$};
  \fill[fmBlue!20] plot[domain=0.75:4, samples=40] (\x,{-1.5+2.0*exp(-(\x-1)/0.6)})
     -- plot[domain=4:0.75, samples=40] (\x,{-1.5-2.0*exp(-(\x-1)/0.6)}) -- cycle;
  \fill[fmGreen!25] plot[domain=0.75:4, samples=40] (\x,{-3.6+1.4*exp(-(\x-1)/0.5)})
     -- plot[domain=4:0.75, samples=40] (\x,{-3.6-1.4*exp(-(\x-1)/0.5)}) -- cycle;
  \fill[fmAmber!35] plot[domain=0.75:4, samples=40] (\x,{-7.5+0.9*exp(-(\x-0.75)/0.2)})
     -- plot[domain=4:0.75, samples=40] (\x,{-7.5-0.9*exp(-(\x-0.75)/0.2)}) -- cycle;
  \node[font=\small, anchor=west] at (4.02,-1.5) {$2p$};
  \node[font=\small, anchor=west] at (4.02,-3.6) {$2s$};
  \node[font=\small, anchor=west] at (4.02,-7.5) {$1s$};
  \draw[fmRed, dashed] (1.25,-9.2) -- (1.25,0.6) node[above, font=\small]{$r_0$};
  \node[font=\scriptsize, fmGray, align=center] at (3.3,0.4) {изолированные\\атомы};
  \node[font=\scriptsize, fmGray, align=center] at (1.0,-8.7) {кристалл};
\end{tikzpicture}
\caption{Расщепление атомных уровней в зоны при сближении атомов
(схема). При равновесном расстоянии $r_0$ уровни $2s$ и $2p$
превратились в широкие перекрывающиеся зоны, а глубокий уровень $1s$
почти не уширился}
\label{fig:levels-bands}
\end{figure}

\paragraph{Двумерная решётка}
В двумерной квадратной решётке с периодом $a$ роль отражающих
«плоскостей» играют семейства параллельных прямых (рис.~\ref{fig:2d-families}).
Пусть волна бежит вдоль оси $x$.
\begin{itemize}
\item Семейство I (прямые $\perp x$): $d=a$, $\theta=90^\circ$, отражение
при $\lambda=2a$, т.\,е. $k_x=\pi/a$.
\item Семейство II (диагонали): $d=a/\sqrt2$, $\theta=45^\circ$,
$\lambda=2\cdot\frac{a}{\sqrt2}\cdot\frac{\sqrt2}{2}=a$, т.\,е. $k_x=2\pi/a$.
\item Семейство III: $d=a/\sqrt5$, $\sin\theta=2/\sqrt5$, $\lambda=4a/5$,
т.\,е. $k_x=5\pi/2a$.
\end{itemize}
Для волны, бегущей под углом $45^\circ$ к оси $x$, семейства I и II дают
разрывы при одном и том же $k=\sqrt2\,\pi/a$. Точки разрыва, таким
образом, зависят от направления движения электрона.

\begin{figure}[ht]
\centering
\begin{tikzpicture}[scale=0.62]
  % семейство I
  \begin{scope}
    \begin{scope}\clip (-0.5,-0.5) rectangle (4.5,4.5);
      \foreach \x in {0,...,4}{\draw[fmBlue, line width=0.7pt] (\x,-1) -- (\x,5);}
    \end{scope}
    \foreach \x in {0,...,4}{\foreach \y in {0,...,4}{\fill[atomLine] (\x,\y) circle (0.09);}}
    \node[font=\small] at (2,-1.1) {I: $d=a$};
  \end{scope}
  % семейство II
  \begin{scope}[shift={(6.2,0)}]
    \begin{scope}\clip (-0.5,-0.5) rectangle (4.5,4.5);
      \foreach \c in {-4,...,8}{\draw[fmGreen!70!black, line width=0.7pt] (\c,0) -- ++(-5,5) (\c,0) -- ++(5,-5);}
    \end{scope}
    \foreach \x in {0,...,4}{\foreach \y in {0,...,4}{\fill[atomLine] (\x,\y) circle (0.09);}}
    \node[font=\small] at (2,-1.1) {II: $d=a/\sqrt2$};
  \end{scope}
  % семейство III: прямые 2x + y = n
  \begin{scope}[shift={(12.4,0)}]
    \begin{scope}\clip (-0.5,-0.5) rectangle (4.5,4.5);
      \foreach \c in {-6,...,14}{\draw[fmRed, line width=0.7pt] ({\c/2},0) -- ++(-3,6) ({\c/2},0) -- ++(3,-6);}
    \end{scope}
    \foreach \x in {0,...,4}{\foreach \y in {0,...,4}{\fill[atomLine] (\x,\y) circle (0.09);}}
    \node[font=\small] at (2,-1.1) {III: $d=a/\sqrt5$};
  \end{scope}
  \draw[arrB] (-0.6,5.0) -- (1.4,5.0) node[right, font=\small, black]{волна вдоль $x$};
\end{tikzpicture}
\caption{Три семейства отражающих прямых в двумерной квадратной решётке.
Для волны, бегущей вдоль $x$, они вызывают разрывы $E(k)$ при
$k_x=\pi/a$, $2\pi/a$ и $5\pi/2a$}
\label{fig:2d-families}
\end{figure}

\paragraph{Зоны Бриллюэна}
В трёхмерной решётке нужно найти все волновые векторы $\vb k$,
удовлетворяющие условию Вульфа--Брэгга. Это удобно сделать с помощью
обратной решётки (п.~1.2.4). Каждому семейству плоскостей $\{hkl\}$
отвечает вектор обратной решётки $\vb G$, перпендикулярный плоскостям,
длиной $G=2\pi/d$\footnote{В кристаллографии (и в оригинале) вектор
обратной решётки определяют без множителя $2\pi$: $H=1/d$. Переход к
пространству волновых векторов~— изменение масштаба в $2\pi$ раз по
всем осям: $\vb G=2\pi\vb H$.}. Из закона Вульфа--Брэгга
$2d\sin\theta=\lambda$ (порядок отражения учтён в $d$) следует
\begin{keyeq}\label{eq:bragg-k}
k\sin\theta=\frac G2 ,
\end{keyeq}
т.\,е. проекция $\vb k$ на направление $\vb G$ равна $G/2$. Значит,
отражение испытывают те и только те $\vb k$, концы которых лежат на
плоскости, \emph{перпендикулярной вектору $\vb G$ и делящей его
пополам} (рис.~\ref{fig:bz-2d},\,\textit{а}). При пересечении такой
\emph{брэгговской плоскости} энергия электрона меняется скачком.

\begin{opredelenie}
\term{Зона Бриллюэна}~— область $\vb k$-пространства, внутри которой
энергия электрона меняется непрерывно, а на границах~— скачком.
\term{Первая зона}~— все точки, которых можно достичь из начала
координат, не пересекая ни одной брэгговской плоскости; \term{$n$-я
зона}~— точки, для достижения которых нужно пересечь ровно $n-1$
плоскостей. Форма и размеры зон полностью определяются кристаллической
структурой и параметром решётки.
\end{opredelenie}

\begin{figure}[ht]
\centering
\begin{subfigure}[b]{0.34\textwidth}\centering
\begin{tikzpicture}[scale=0.72]
  \draw[-{Stealth[length=2mm]}, fmGray] (-0.5,0) -- (3.6,0) node[right, font=\small]{$k_x$};
  \draw[-{Stealth[length=2mm]}, fmGray] (0,-0.5) -- (0,3.4);
  \draw[fmRed, line width=0.9pt] (1.5,-0.6) -- (1.5,3.2) node[above, font=\small]{MM};
  \draw[fmRed, line width=0.9pt] (-1.5,-0.6) -- (-1.5,3.2) node[above, font=\small]{NN};
  \draw[-{Stealth[length=2.4mm]}, line width=1.1pt] (0,0) -- (3,0);
  \node[font=\small, below] at (2.6,0) {$\vb G$};
  \draw[arr, fmBlue] (0,0) -- (1.5,2.2) node[midway, left, font=\small]{$\vb k$};
  \draw[dimline] (0,-0.35) -- node[below, font=\scriptsize]{$G/2$} (1.5,-0.35);
  \fill (0,0) circle (1.5pt);
\end{tikzpicture}
\caption{}
\end{subfigure}\hfill
\begin{subfigure}[b]{0.6\textwidth}\centering
\begin{tikzpicture}[scale=0.62]
  % зона 3: восемь треугольников
  \foreach \sx in {-1,1}{\foreach \sy in {-1,1}{
    \fill[fmAmber!35] (2*\sx,0) -- (\sx,\sy) -- (2*\sx,\sy) -- cycle;
    \fill[fmAmber!35] (0,2*\sy) -- (\sx,\sy) -- (\sx,2*\sy) -- cycle;}}
  % зона 2: четыре треугольника
  \fill[fmGreen!30] (1,-1) -- (2,0) -- (1,1) -- cycle;
  \fill[fmGreen!30] (-1,-1) -- (-2,0) -- (-1,1) -- cycle;
  \fill[fmGreen!30] (-1,1) -- (0,2) -- (1,1) -- cycle;
  \fill[fmGreen!30] (-1,-1) -- (0,-2) -- (1,-1) -- cycle;
  % зона 1
  \fill[fmBlue!30] (-1,-1) rectangle (1,1);
  % брэгговские прямые
  \begin{scope}\clip (-2.6,-2.6) rectangle (2.6,2.6);
    \foreach \c in {-2,-1,1,2}{\draw[atomLine, line width=0.4pt] (\c,-3) -- (\c,3) (-3,\c) -- (3,\c);}
    \draw[atomLine, line width=0.4pt] (-3,1) -- (1,-3) (-1,3) -- (3,-1) (-3,-1) -- (1,3) (-1,-3) -- (3,1);
  \end{scope}
  \node[font=\small] at (0,0) {1};
  \node[font=\small] at (1.4,0) {2}; \node[font=\small] at (-1.4,0) {2};
  \node[font=\small] at (0,1.4) {2}; \node[font=\small] at (0,-1.4) {2};
  \foreach \sx in {-1,1}{\foreach \sy in {-1,1}{
    \node[font=\scriptsize] at (1.72*\sx,0.62*\sy) {3}; \node[font=\scriptsize] at (0.62*\sx,1.72*\sy) {3};}}
  \draw[-{Stealth[length=2mm]}, fmGray] (-2.8,0) -- (3.0,0) node[right, font=\small]{$k_x$};
  \draw[-{Stealth[length=2mm]}, fmGray] (0,-2.8) -- (0,3.0) node[above, font=\small]{$k_y$};
  \node[font=\scriptsize, below right] at (1,0) {$\frac\pi a$};
  \node[font=\scriptsize, below right] at (2,0) {$\frac{2\pi}a$};
\end{tikzpicture}
\caption{}
\end{subfigure}
\caption{\textit{а}~— условие отражения в $\vb k$-пространстве: концы
векторов $\vb k$, испытывающих брэгговское отражение, лежат на
плоскостях MM и NN, делящих вектор обратной решётки $\vb G$ пополам;
\textit{б}~— первые три зоны Бриллюэна квадратной решётки (тонкие
линии~— брэгговские прямые). Площади всех зон одинаковы}
\label{fig:bz-2d}
\end{figure}

\paragraph{Зоны Бриллюэна кубических решёток}
Из \eqref{eq:bragg-k} видно, что ближе всего к началу координат лежат
брэгговские плоскости семейства с \emph{наибольшим} межплоскостным
расстоянием $d_\text{max}$: они удалены на $G/2=\pi/d_\text{max}$.
Именно они в первую очередь ограничивают первую зону
(рис.~\ref{fig:bz-3d}).
\begin{itemize}
\item \emph{Простая кубическая} решётка (чистые металлы её не имеют,
кроме $\alpha$-Po, но она лежит в основе упорядоченной $\beta$-латуни и
CsCl). Наибольшее $d$ у плоскостей $\{100\}$: $d=a$. Шесть брэгговских
плоскостей на расстоянии $\pi/a$ образуют \emph{куб} с ребром $2\pi/a$.
Вторую зону ограничивают плоскости $\{110\}$ ($d=a/\sqrt2$), удалённые
на $\sqrt2\,\pi/a$\ispr{В оригинале (с.~391) вторая зона связана с
семейством $\{100\}$ с $d=a/\sqrt2$; такое межплоскостное расстояние
имеют плоскости $\{110\}$.}~— двумерное сечение показано на
рис.~\ref{fig:bz-2d},\,\textit{б}.
\item \emph{ГЦК-решётка} (Cu, Ag, Au). Наибольшее $d=a/\sqrt3$ у
плоскостей $\{111\}$: восемь плоскостей на расстоянии $\sqrt3\,\pi/a$
образуют октаэдр. Его вершины срезают шесть плоскостей $\{200\}$
($d=a/2$), удалённые на $2\pi/a$\ispr{В оригинале это семейство
названо $\{100\}$. В ГЦК-решётке отражения от $\{100\}$ погашаются, и
первое отражение этого типа даёт семейство $\{200\}$ с $d=a/2$~—
именно ему соответствует расстояние $2\pi/a$.}. Первая зона~—
\term{усечённый октаэдр} с восемью шестиугольными и шестью квадратными
гранями\ispr{В оригинале (с.~392) эта фигура названа кубооктаэдром.
Кубооктаэдр имеет треугольные и квадратные грани; первая зона ГЦК-решётки~—
усечённый октаэдр (шестиугольные и квадратные грани).}.
\item \emph{ОЦК-решётка} (щелочные металлы, Fe, W, Mo). Наибольшее
$d=a/\sqrt2$ у плоскостей $\{110\}$; двенадцать плоскостей на расстоянии
$\sqrt2\,\pi/a$ образуют \term{ромбический додекаэдр}.
\end{itemize}

\begin{figure}[ht]
\centering
\begin{subfigure}[b]{0.3\textwidth}\centering
\begin{tikzpicture}[scale=1.55]
  % сгенерировано tools/gen_bands.py — не править вручную
\fill[fmBlue!17] (-1.014,0.617) -- (-1.014,-0.809) -- (0.310,-1.027) -- (0.310,0.400) -- cycle;
\fill[fmBlue!11] (0.310,-1.027) -- (1.014,-0.617) -- (1.014,0.809) -- (0.310,0.400) -- cycle;
\fill[fmBlue!20] (1.014,0.809) -- (-0.310,1.027) -- (-1.014,0.617) -- (0.310,0.400) -- cycle;
\draw[bzfront] (-1.014,-0.809) -- (-1.014,0.617);
\draw[bzfront] (-1.014,-0.809) -- (0.310,-1.027);
\draw[bzfront] (0.310,-1.027) -- (0.310,0.400);
\draw[bzfront] (-1.014,0.617) -- (0.310,0.400);
\draw[bzback] (-0.310,-0.400) -- (1.014,-0.617);
\draw[bzback] (-0.310,-0.400) -- (-0.310,1.027);
\draw[bzfront] (-0.310,1.027) -- (1.014,0.809);
\draw[bzfront] (1.014,-0.617) -- (1.014,0.809);
\draw[bzfront] (1.014,-0.617) -- (0.310,-1.027);
\draw[bzfront] (1.014,0.809) -- (0.310,0.400);
\draw[bzback] (-0.310,-0.400) -- (-1.014,-0.809);
\draw[bzfront] (-0.310,1.027) -- (-1.014,0.617);
\coordinate (bzG) at (0.000,0.000);
\coordinate (bzX) at (-0.352,-0.205);
\coordinate (bzR) at (0.310,0.400);
\coordinate (bzM) at (0.310,-0.313);

  \draw[fmRed, line width=0.6pt] (bzG) -- (bzX);
  \node[bzpt] at (bzG) {}; \node[bzpt] at (bzX) {}; \node[bzpt] at (bzR) {};
  \node[font=\scriptsize, above left=-1pt] at (bzG) {$\Gamma$};
  \node[font=\scriptsize, below right=-1pt] at (bzX) {$X$};
  \node[font=\scriptsize, above right=-1pt] at (bzR) {$R$};
\end{tikzpicture}
\caption{простая кубическая}
\end{subfigure}\hfill
\begin{subfigure}[b]{0.32\textwidth}\centering
\begin{tikzpicture}[scale=1.55]
  % сгенерировано tools/gen_bands.py — не править вручную
\fill[fmBlue!17] (-0.469,0.203) -- (-0.911,-0.200) -- (-0.469,-0.748) -- (-0.028,-0.345) -- cycle;
\fill[fmBlue!11] (0.883,-0.621) -- (1.118,-0.009) -- (0.883,0.330) -- (0.648,-0.281) -- cycle;
\fill[fmBlue!20] (0.441,0.879) -- (0.235,1.087) -- (-0.441,1.024) -- (-0.235,0.815) -- cycle;
\fill[fmBlue!15] (-0.911,-0.200) -- (-1.118,0.009) -- (-0.883,-0.330) -- (-0.441,-0.879) -- (-0.235,-1.087) -- (-0.469,-0.748) -- cycle;
\fill[fmBlue!21] (-0.235,0.815) -- (-0.441,1.024) -- (-0.883,0.621) -- (-1.118,0.009) -- (-0.911,-0.200) -- (-0.469,0.203) -- cycle;
\fill[fmBlue!10] (-0.469,-0.748) -- (-0.235,-1.087) -- (0.441,-1.024) -- (0.883,-0.621) -- (0.648,-0.281) -- (-0.028,-0.345) -- cycle;
\fill[fmBlue!17] (0.648,-0.281) -- (0.883,0.330) -- (0.441,0.879) -- (-0.235,0.815) -- (-0.469,0.203) -- (-0.028,-0.345) -- cycle;
\draw[bzfront] (-0.911,-0.200) -- (-0.469,0.203);
\draw[bzfront] (-0.911,-0.200) -- (-0.469,-0.748);
\draw[bzfront] (-0.469,-0.748) -- (-0.028,-0.345);
\draw[bzfront] (-0.469,0.203) -- (-0.028,-0.345);
\draw[bzback] (0.028,0.345) -- (0.469,-0.203);
\draw[bzback] (0.028,0.345) -- (0.469,0.748);
\draw[bzback] (0.469,0.748) -- (0.911,0.200);
\draw[bzback] (0.469,-0.203) -- (0.911,0.200);
\draw[bzfront] (1.118,-0.009) -- (0.883,-0.621);
\draw[bzfront] (1.118,-0.009) -- (0.883,0.330);
\draw[bzfront] (0.883,0.330) -- (0.648,-0.281);
\draw[bzfront] (0.883,-0.621) -- (0.648,-0.281);
\draw[bzback] (-0.648,0.281) -- (-0.883,0.621);
\draw[bzback] (-0.648,0.281) -- (-0.883,-0.330);
\draw[bzfront] (-0.883,-0.330) -- (-1.118,0.009);
\draw[bzfront] (-0.883,0.621) -- (-1.118,0.009);
\draw[bzfront] (0.235,1.087) -- (0.441,0.879);
\draw[bzfront] (0.235,1.087) -- (-0.441,1.024);
\draw[bzfront] (-0.441,1.024) -- (-0.235,0.815);
\draw[bzfront] (0.441,0.879) -- (-0.235,0.815);
\draw[bzback] (0.235,-0.815) -- (-0.441,-0.879);
\draw[bzback] (0.235,-0.815) -- (0.441,-1.024);
\draw[bzfront] (0.441,-1.024) -- (-0.235,-1.087);
\draw[bzfront] (-0.441,-0.879) -- (-0.235,-1.087);
\draw[bzback] (0.028,0.345) -- (-0.648,0.281);
\draw[bzback] (0.469,-0.203) -- (0.235,-0.815);
\draw[bzfront] (-0.883,-0.330) -- (-0.441,-0.879);
\draw[bzback] (0.469,0.748) -- (0.235,1.087);
\draw[bzfront] (-0.883,0.621) -- (-0.441,1.024);
\draw[bzback] (0.911,0.200) -- (1.118,-0.009);
\draw[bzfront] (0.441,-1.024) -- (0.883,-0.621);
\draw[bzfront] (0.441,0.879) -- (0.883,0.330);
\draw[bzfront] (-1.118,0.009) -- (-0.911,-0.200);
\draw[bzfront] (-0.235,-1.087) -- (-0.469,-0.748);
\draw[bzfront] (-0.235,0.815) -- (-0.469,0.203);
\draw[bzfront] (0.648,-0.281) -- (-0.028,-0.345);
\coordinate (bzG) at (0.000,0.000);
\coordinate (bzX) at (-0.469,-0.273);
\coordinate (bzL) at (0.207,0.267);
\coordinate (bzW) at (-0.028,-0.345);

  \draw[fmRed, line width=0.6pt] (bzG) -- (bzX) (bzG) -- (bzL);
  \foreach \p in {G,X,L,W}{\node[bzpt] at (bz\p) {};}
  \node[font=\scriptsize, above left=-1pt] at (bzG) {$\Gamma$};
  \node[font=\scriptsize, below=1pt] at (bzX) {$X$};
  \node[font=\scriptsize, above right=-1pt] at (bzL) {$L$};
  \node[font=\scriptsize, right=1pt] at (bzW) {$W$};
\end{tikzpicture}
\caption{ГЦК}
\end{subfigure}\hfill
\begin{subfigure}[b]{0.32\textwidth}\centering
\begin{tikzpicture}[scale=1.55]
  % сгенерировано tools/gen_bands.py — не править вручную
\fill[fmBlue!18] (-0.676,-0.539) -- (-0.469,-0.273) -- (-0.676,0.412) -- (-0.883,0.145) -- cycle;
\fill[fmBlue!14] (0.207,0.267) -- (-0.469,-0.273) -- (0.207,-0.684) -- (0.883,-0.145) -- cycle;
\fill[fmBlue!12] (0.207,-0.684) -- (-0.469,-0.273) -- (-0.676,-0.539) -- (0.000,-0.951) -- cycle;
\fill[fmBlue!20] (-0.676,0.412) -- (-0.469,-0.273) -- (0.207,0.267) -- (0.000,0.951) -- cycle;
\fill[fmBlue!8] (0.676,-0.412) -- (0.883,-0.145) -- (0.207,-0.684) -- (0.000,-0.951) -- cycle;
\fill[fmBlue!16] (0.207,0.267) -- (0.883,-0.145) -- (0.676,0.539) -- (0.000,0.951) -- cycle;
\draw[bzback] (0.469,0.273) -- (-0.207,0.684);
\draw[bzback] (0.469,0.273) -- (-0.207,-0.267);
\draw[bzback] (-0.883,0.145) -- (-0.207,-0.267);
\draw[bzback] (-0.883,0.145) -- (-0.207,0.684);
\draw[bzback] (0.469,0.273) -- (0.676,-0.412);
\draw[bzback] (0.469,0.273) -- (0.676,0.539);
\draw[bzfront] (0.883,-0.145) -- (0.676,0.539);
\draw[bzfront] (0.883,-0.145) -- (0.676,-0.412);
\draw[bzfront] (-0.469,-0.273) -- (-0.676,-0.539);
\draw[bzfront] (-0.469,-0.273) -- (-0.676,0.412);
\draw[bzfront] (-0.883,0.145) -- (-0.676,0.412);
\draw[bzfront] (-0.883,0.145) -- (-0.676,-0.539);
\draw[bzfront] (-0.469,-0.273) -- (0.207,0.267);
\draw[bzfront] (-0.469,-0.273) -- (0.207,-0.684);
\draw[bzfront] (0.883,-0.145) -- (0.207,-0.684);
\draw[bzfront] (0.883,-0.145) -- (0.207,0.267);
\draw[bzfront] (0.000,-0.951) -- (0.676,-0.412);
\draw[bzback] (0.000,-0.951) -- (-0.207,-0.267);
\draw[bzback] (0.000,0.951) -- (-0.207,0.684);
\draw[bzfront] (0.000,0.951) -- (0.676,0.539);
\draw[bzfront] (0.000,-0.951) -- (-0.676,-0.539);
\draw[bzfront] (0.000,-0.951) -- (0.207,-0.684);
\draw[bzfront] (0.000,0.951) -- (0.207,0.267);
\draw[bzfront] (0.000,0.951) -- (-0.676,0.412);
\coordinate (bzG) at (0.000,0.000);
\coordinate (bzH) at (-0.469,-0.273);
\coordinate (bzN) at (0.207,-0.209);
\coordinate (bzP) at (0.207,0.267);

  \draw[fmRed, line width=0.6pt] (bzG) -- (bzH) (bzG) -- (bzN);
  \foreach \p in {G,H,N,P}{\node[bzpt] at (bz\p) {};}
  \node[font=\scriptsize, above left=-1pt] at (bzG) {$\Gamma$};
  \node[font=\scriptsize, below=1pt] at (bzH) {$H$};
  \node[font=\scriptsize, right=1pt] at (bzN) {$N$};
  \node[font=\scriptsize, above right=-1pt] at (bzP) {$P$};
\end{tikzpicture}
\caption{ОЦК}
\end{subfigure}
\caption{Первые зоны Бриллюэна кубических решёток (расчёт; невидимые рёбра~—
пунктиром): \textit{а}~— куб; \textit{б}~— усечённый октаэдр;
\textit{в}~— ромбический додекаэдр. Красным отмечены точки высокой
симметрии: $\Gamma$~— центр зоны; $X$, $L$, $H$, $N$~— центры граней;
$W$, $R$, $P$~— вершины}
\label{fig:bz-3d}
\end{figure}

\paragraph{Энергии в точках зоны}
Оценим энергию свободного электрона в характерных точках зоны. Для меди
($a=0{,}361$~нм) в центре квадратной грани (точка $X$)
\[
k_X=\frac{2\pi}a=17{,}4~\text{нм}^{-1},\qquad
E_X=\frac{\hbar^2k_X^2}{2m}=0{,}0381\cdot17{,}4^2\approx11{,}5~\text{эВ}.
\]
Так же рассчитаны остальные значения в табл.~\ref{tab:bz-points}. В
последнем столбце~— радиус сферы, объём которой равен \emph{половине}
объёма зоны, и энергия на её поверхности. Сравните её с энергией Ферми
тех же металлов в табл.~\ref{tab:EF}: совпадение почти полное. Почему
именно половина зоны~— выясним сейчас.

\begin{table}[ht]
\centering
\caption{Энергии свободного электрона в точках зон Бриллюэна (эВ).
В оригинале точки $L$, $X$, $W$ обозначены $A$, $B$, $C$, а точки $N$,
$H$~— $A$, $B$}
\label{tab:bz-points}
\small
\begin{tabular}{@{}lcccc@{\qquad}lccc@{}}
\toprule
\multicolumn{5}{@{}l}{ГЦК} & \multicolumn{4}{l}{ОЦК} \\
\cmidrule(r){1-5}\cmidrule(l){6-9}
Точка & $L$ & $X$ & $W$ & $\frac12$ зоны & Точка & $N$ & $H$ & $\frac12$ зоны \\
$k$, ед. $2\pi/a$ & 0,866 & 1 & 1,118 & 0,782 & $k$, ед. $2\pi/a$ & 0,707 & 1 & 0,620 \\
\midrule
Cu & 8,64 & 11,52 & 14,40 & 7,04 & Li & 6,12 & 12,25 & 4,71 \\
Ag & 6,75 & 9,00 & 11,27 & 5,51 & Na & 4,05 & 8,11 & 3,12 \\
Au & 6,79 & 9,06 & 11,32 & 5,53 & K & 2,77 & 5,55 & 2,13 \\
\bottomrule
\end{tabular}
\end{table}
\noindent Энергия серебра в точке $X$ исправлена\ispr{В табл.~3.13
оригинала для серебра в точке $B$ ($X$) указано 9,2~эВ. Но энергия
пропорциональна $k^2$, и из значения в точке $A$ ($L$) следует
$6{,}746/0{,}866^2=9{,}0$~эВ; так же получается и из параметра решётки
$a=0{,}409$~нм.}.

\paragraph{Сколько электронов вмещает зона}
Разделим объём зоны на объём одного состояния. Для простой кубической
решётки объём первой зоны $(2\pi/a)^3$, объём одного состояния
$8\pi^3/V$; их отношение
\[
\frac{(2\pi/a)^3}{8\pi^3/V}=\frac V{a^3}=N
\]
равно числу элементарных ячеек, т.\,е. числу атомов. В каждом состоянии~—
два электрона с противоположными спинами.
\begin{glavnoe}
Зона Бриллюэна простой кубической, ГЦК- и ОЦК-решётки вмещает \vazhno{по
два электрона на атом}. У одновалентных металлов (щелочные, Cu, Ag, Au)
электроны заполняют поэтому лишь \emph{половину} первой зоны~— сферу,
равную по объёму половине зоны. Отсюда и совпадение в
табл.~\ref{tab:bz-points}.
\end{glavnoe}
Для ГПУ-решётки зоны не содержат целого числа состояний на атом:
число зависит от отношения осей, и, например, у цинка ($c/a=1{,}86$)
первая зона вмещает лишь 1,792 электрона на атом.

\begin{primer}[Когда сфера Ферми касается границы зоны]
Сколько электронов на атом поместится в сфере Ферми, \emph{вписанной}
в первую зону ОЦК-решётки?
\begin{shagi}
\item Радиус вписанной сферы~— расстояние до грани $\{110\}$:
$r=\sqrt2\,\pi/a=0{,}707\cdot(2\pi/a)$.
\item Объём сферы $\frac43\pi r^3=1{,}481\,(2\pi/a)^3$; объём зоны
$2\,(2\pi/a)^3$; отношение $0{,}74$.
\item Полная зона вмещает 2 электрона на атом, значит вписанная сфера~—
$2\cdot0{,}74=1{,}48$ электрона на атом.
\end{shagi}
Для ГЦК-решётки (вписанная сфера касается граней $\{111\}$,
$r=\sqrt3\,\pi/a$; объём зоны $4(2\pi/a)^3$) отношение объёмов $0{,}68$,
и электронная концентрация равна $1{,}36$ электрона на атом. Эти числа
понадобятся в п.~3.3 при объяснении электронных соединений.
\end{primer}

\paragraph{Как заполняется зона}
Проследим, как меняется изоэнергетическая поверхность при добавлении
электронов, на примере двумерной квадратной решётки
(рис.~\ref{fig:contours}). Минимум энергии~— в центре зоны.
\begin{shagi}
\item При малых $k$ электроны почти свободны: $E\approx\hbar^2k^2/2m$,
изоэнергетические линии~— окружности.
\item Вблизи границ $E(k)$ отстаёт от параболы, причём раньше всего~— в
направлении~1, где граница ближе (до середины стороны $C$ расстояние
$\pi/a$, а до угла~— $\sqrt2\,\pi/a$). Поэтому при одной и той же энергии
$k$ в направлении~1 больше: линии «вытягиваются» к серединам сторон.
\item Коснуться границы линия не может: вблизи точек $C$ она
перестраивается~— разрывается и подходит к границе под прямым углом.
Замкнутая (односвязная) линия становится многосвязной.
\item Дальше остаются незаполненными только углы зоны.
\end{shagi}

\begin{figure}[ht]
\centering
\begin{tikzpicture}[scale=2.0,
  c1/.style={fmBlue, line width=0.8pt}, c2/.style={fmBlue!80!fmRed, line width=0.8pt},
  c3/.style={fmBlue!65!fmRed, line width=0.8pt}, c4/.style={fmBlue!50!fmRed, line width=0.8pt},
  c5/.style={fmBlue!35!fmRed, line width=0.8pt}, c6/.style={fmRed!80!black, line width=1.2pt},
  c7/.style={fmRed!70!fmAmber, line width=0.8pt}, c8/.style={fmAmber, line width=0.8pt}]
  \fill[fmBlue!5] (-1,-1) rectangle (1,1);
  % сгенерировано tools/gen_bands.py — не править вручную
\draw[c1] plot[smooth] coordinates {(-0.230,0.000) (-0.230,0.008) (-0.229,0.016) (-0.229,0.023) (-0.228,0.031) (-0.226,0.039) (-0.225,0.046) (-0.223,0.054) (-0.221,0.062) (-0.219,0.069) (-0.216,0.076) (-0.213,0.084) (-0.210,0.091) (-0.206,0.098) (-0.203,0.105) (-0.199,0.112) (-0.194,0.119) (-0.190,0.126) (-0.185,0.133) (-0.180,0.139) (-0.175,0.145) (-0.170,0.152) (-0.164,0.158) (-0.158,0.163) (-0.152,0.169) (-0.146,0.174) (-0.140,0.180) (-0.133,0.185) (-0.127,0.189) (-0.120,0.194) (-0.113,0.198) (-0.106,0.202) (-0.098,0.206) (-0.091,0.210) (-0.083,0.213) (-0.076,0.216) (-0.068,0.219) (-0.060,0.221) (-0.052,0.223) (-0.044,0.225) (-0.036,0.227) (-0.028,0.228) (-0.020,0.229) (-0.012,0.230) (-0.004,0.230) (0.004,0.230) (0.012,0.230) (0.020,0.229) (0.028,0.228) (0.036,0.227) (0.044,0.225) (0.052,0.223) (0.060,0.221) (0.068,0.219) (0.076,0.216) (0.083,0.213) (0.091,0.210) (0.098,0.206) (0.106,0.202) (0.113,0.198) (0.120,0.194) (0.127,0.189) (0.133,0.185) (0.140,0.180) (0.146,0.174) (0.152,0.169) (0.158,0.163) (0.164,0.158) (0.170,0.152) (0.175,0.145) (0.180,0.139) (0.185,0.133) (0.190,0.126) (0.194,0.119) (0.199,0.112) (0.203,0.105) (0.206,0.098) (0.210,0.091) (0.213,0.084) (0.216,0.076) (0.219,0.069) (0.221,0.062) (0.223,0.054) (0.225,0.046) (0.226,0.039) (0.228,0.031) (0.229,0.023) (0.229,0.016) (0.230,0.008) (0.230,0.000)};
\draw[c1] plot[smooth] coordinates {(-0.230,-0.000) (-0.230,-0.008) (-0.229,-0.016) (-0.229,-0.023) (-0.228,-0.031) (-0.226,-0.039) (-0.225,-0.046) (-0.223,-0.054) (-0.221,-0.062) (-0.219,-0.069) (-0.216,-0.076) (-0.213,-0.084) (-0.210,-0.091) (-0.206,-0.098) (-0.203,-0.105) (-0.199,-0.112) (-0.194,-0.119) (-0.190,-0.126) (-0.185,-0.133) (-0.180,-0.139) (-0.175,-0.145) (-0.170,-0.152) (-0.164,-0.158) (-0.158,-0.163) (-0.152,-0.169) (-0.146,-0.174) (-0.140,-0.180) (-0.133,-0.185) (-0.127,-0.189) (-0.120,-0.194) (-0.113,-0.198) (-0.106,-0.202) (-0.098,-0.206) (-0.091,-0.210) (-0.083,-0.213) (-0.076,-0.216) (-0.068,-0.219) (-0.060,-0.221) (-0.052,-0.223) (-0.044,-0.225) (-0.036,-0.227) (-0.028,-0.228) (-0.020,-0.229) (-0.012,-0.230) (-0.004,-0.230) (0.004,-0.230) (0.012,-0.230) (0.020,-0.229) (0.028,-0.228) (0.036,-0.227) (0.044,-0.225) (0.052,-0.223) (0.060,-0.221) (0.068,-0.219) (0.076,-0.216) (0.083,-0.213) (0.091,-0.210) (0.098,-0.206) (0.106,-0.202) (0.113,-0.198) (0.120,-0.194) (0.127,-0.189) (0.133,-0.185) (0.140,-0.180) (0.146,-0.174) (0.152,-0.169) (0.158,-0.163) (0.164,-0.158) (0.170,-0.152) (0.175,-0.145) (0.180,-0.139) (0.185,-0.133) (0.190,-0.126) (0.194,-0.119) (0.199,-0.112) (0.203,-0.105) (0.206,-0.098) (0.210,-0.091) (0.213,-0.084) (0.216,-0.076) (0.219,-0.069) (0.221,-0.062) (0.223,-0.054) (0.225,-0.046) (0.226,-0.039) (0.228,-0.031) (0.229,-0.023) (0.229,-0.016) (0.230,-0.008) (0.230,-0.000)};
\draw[c2] plot[smooth] coordinates {(-0.369,0.000) (-0.369,0.012) (-0.368,0.023) (-0.367,0.035) (-0.365,0.046) (-0.363,0.058) (-0.361,0.069) (-0.358,0.081) (-0.354,0.092) (-0.351,0.104) (-0.346,0.115) (-0.342,0.126) (-0.336,0.137) (-0.331,0.148) (-0.325,0.159) (-0.318,0.170) (-0.312,0.181) (-0.305,0.192) (-0.297,0.202) (-0.289,0.213) (-0.281,0.223) (-0.272,0.233) (-0.263,0.243) (-0.254,0.252) (-0.244,0.262) (-0.234,0.271) (-0.224,0.280) (-0.214,0.288) (-0.203,0.297) (-0.192,0.304) (-0.181,0.312) (-0.169,0.319) (-0.158,0.326) (-0.146,0.332) (-0.134,0.338) (-0.121,0.344) (-0.109,0.349) (-0.097,0.353) (-0.084,0.357) (-0.071,0.360) (-0.058,0.363) (-0.045,0.365) (-0.033,0.367) (-0.020,0.368) (-0.007,0.369) (0.007,0.369) (0.020,0.368) (0.033,0.367) (0.045,0.365) (0.058,0.363) (0.071,0.360) (0.084,0.357) (0.097,0.353) (0.109,0.349) (0.121,0.344) (0.134,0.338) (0.146,0.332) (0.158,0.326) (0.169,0.319) (0.181,0.312) (0.192,0.304) (0.203,0.297) (0.214,0.288) (0.224,0.280) (0.234,0.271) (0.244,0.262) (0.254,0.252) (0.263,0.243) (0.272,0.233) (0.281,0.223) (0.289,0.213) (0.297,0.202) (0.305,0.192) (0.312,0.181) (0.318,0.170) (0.325,0.159) (0.331,0.148) (0.336,0.137) (0.342,0.126) (0.346,0.115) (0.351,0.104) (0.354,0.092) (0.358,0.081) (0.361,0.069) (0.363,0.058) (0.365,0.046) (0.367,0.035) (0.368,0.023) (0.369,0.012) (0.369,0.000)};
\draw[c2] plot[smooth] coordinates {(-0.369,-0.000) (-0.369,-0.012) (-0.368,-0.023) (-0.367,-0.035) (-0.365,-0.046) (-0.363,-0.058) (-0.361,-0.069) (-0.358,-0.081) (-0.354,-0.092) (-0.351,-0.104) (-0.346,-0.115) (-0.342,-0.126) (-0.336,-0.137) (-0.331,-0.148) (-0.325,-0.159) (-0.318,-0.170) (-0.312,-0.181) (-0.305,-0.192) (-0.297,-0.202) (-0.289,-0.213) (-0.281,-0.223) (-0.272,-0.233) (-0.263,-0.243) (-0.254,-0.252) (-0.244,-0.262) (-0.234,-0.271) (-0.224,-0.280) (-0.214,-0.288) (-0.203,-0.297) (-0.192,-0.304) (-0.181,-0.312) (-0.169,-0.319) (-0.158,-0.326) (-0.146,-0.332) (-0.134,-0.338) (-0.121,-0.344) (-0.109,-0.349) (-0.097,-0.353) (-0.084,-0.357) (-0.071,-0.360) (-0.058,-0.363) (-0.045,-0.365) (-0.033,-0.367) (-0.020,-0.368) (-0.007,-0.369) (0.007,-0.369) (0.020,-0.368) (0.033,-0.367) (0.045,-0.365) (0.058,-0.363) (0.071,-0.360) (0.084,-0.357) (0.097,-0.353) (0.109,-0.349) (0.121,-0.344) (0.134,-0.338) (0.146,-0.332) (0.158,-0.326) (0.169,-0.319) (0.181,-0.312) (0.192,-0.304) (0.203,-0.297) (0.214,-0.288) (0.224,-0.280) (0.234,-0.271) (0.244,-0.262) (0.254,-0.252) (0.263,-0.243) (0.272,-0.233) (0.281,-0.223) (0.289,-0.213) (0.297,-0.202) (0.305,-0.192) (0.312,-0.181) (0.318,-0.170) (0.325,-0.159) (0.331,-0.148) (0.336,-0.137) (0.342,-0.126) (0.346,-0.115) (0.351,-0.104) (0.354,-0.092) (0.358,-0.081) (0.361,-0.069) (0.363,-0.058) (0.365,-0.046) (0.367,-0.035) (0.368,-0.023) (0.369,-0.012) (0.369,-0.000)};
\draw[c3] plot[smooth] coordinates {(-0.500,0.000) (-0.500,0.014) (-0.499,0.028) (-0.497,0.042) (-0.495,0.056) (-0.492,0.070) (-0.489,0.085) (-0.485,0.099) (-0.480,0.113) (-0.475,0.127) (-0.469,0.141) (-0.463,0.155) (-0.456,0.169) (-0.448,0.184) (-0.440,0.198) (-0.432,0.212) (-0.422,0.226) (-0.413,0.240) (-0.402,0.254) (-0.392,0.268) (-0.380,0.282) (-0.369,0.295) (-0.357,0.309) (-0.344,0.322) (-0.331,0.336) (-0.318,0.349) (-0.304,0.361) (-0.290,0.374) (-0.275,0.386) (-0.260,0.398) (-0.245,0.409) (-0.229,0.420) (-0.214,0.430) (-0.197,0.440) (-0.181,0.450) (-0.165,0.458) (-0.148,0.466) (-0.131,0.473) (-0.114,0.480) (-0.096,0.485) (-0.079,0.490) (-0.062,0.494) (-0.044,0.497) (-0.026,0.499) (-0.009,0.500) (0.009,0.500) (0.026,0.499) (0.044,0.497) (0.062,0.494) (0.079,0.490) (0.096,0.485) (0.114,0.480) (0.131,0.473) (0.148,0.466) (0.165,0.458) (0.181,0.450) (0.197,0.440) (0.214,0.430) (0.229,0.420) (0.245,0.409) (0.260,0.398) (0.275,0.386) (0.290,0.374) (0.304,0.361) (0.318,0.349) (0.331,0.336) (0.344,0.322) (0.357,0.309) (0.369,0.295) (0.380,0.282) (0.392,0.268) (0.402,0.254) (0.413,0.240) (0.422,0.226) (0.432,0.212) (0.440,0.198) (0.448,0.184) (0.456,0.169) (0.463,0.155) (0.469,0.141) (0.475,0.127) (0.480,0.113) (0.485,0.099) (0.489,0.085) (0.492,0.070) (0.495,0.056) (0.497,0.042) (0.499,0.028) (0.500,0.014) (0.500,0.000)};
\draw[c3] plot[smooth] coordinates {(-0.500,-0.000) (-0.500,-0.014) (-0.499,-0.028) (-0.497,-0.042) (-0.495,-0.056) (-0.492,-0.070) (-0.489,-0.085) (-0.485,-0.099) (-0.480,-0.113) (-0.475,-0.127) (-0.469,-0.141) (-0.463,-0.155) (-0.456,-0.169) (-0.448,-0.184) (-0.440,-0.198) (-0.432,-0.212) (-0.422,-0.226) (-0.413,-0.240) (-0.402,-0.254) (-0.392,-0.268) (-0.380,-0.282) (-0.369,-0.295) (-0.357,-0.309) (-0.344,-0.322) (-0.331,-0.336) (-0.318,-0.349) (-0.304,-0.361) (-0.290,-0.374) (-0.275,-0.386) (-0.260,-0.398) (-0.245,-0.409) (-0.229,-0.420) (-0.214,-0.430) (-0.197,-0.440) (-0.181,-0.450) (-0.165,-0.458) (-0.148,-0.466) (-0.131,-0.473) (-0.114,-0.480) (-0.096,-0.485) (-0.079,-0.490) (-0.062,-0.494) (-0.044,-0.497) (-0.026,-0.499) (-0.009,-0.500) (0.009,-0.500) (0.026,-0.499) (0.044,-0.497) (0.062,-0.494) (0.079,-0.490) (0.096,-0.485) (0.114,-0.480) (0.131,-0.473) (0.148,-0.466) (0.165,-0.458) (0.181,-0.450) (0.197,-0.440) (0.214,-0.430) (0.229,-0.420) (0.245,-0.409) (0.260,-0.398) (0.275,-0.386) (0.290,-0.374) (0.304,-0.361) (0.318,-0.349) (0.331,-0.336) (0.344,-0.322) (0.357,-0.309) (0.369,-0.295) (0.380,-0.282) (0.392,-0.268) (0.402,-0.254) (0.413,-0.240) (0.422,-0.226) (0.432,-0.212) (0.440,-0.198) (0.448,-0.184) (0.456,-0.169) (0.463,-0.155) (0.469,-0.141) (0.475,-0.127) (0.480,-0.113) (0.485,-0.099) (0.489,-0.085) (0.492,-0.070) (0.495,-0.056) (0.497,-0.042) (0.499,-0.028) (0.500,-0.014) (0.500,-0.000)};
\draw[c4] plot[smooth] coordinates {(-0.631,0.000) (-0.631,0.015) (-0.629,0.030) (-0.627,0.046) (-0.625,0.061) (-0.621,0.076) (-0.617,0.091) (-0.612,0.107) (-0.606,0.122) (-0.599,0.138) (-0.592,0.154) (-0.584,0.170) (-0.575,0.186) (-0.566,0.202) (-0.555,0.219) (-0.545,0.235) (-0.533,0.252) (-0.521,0.268) (-0.508,0.285) (-0.494,0.302) (-0.480,0.319) (-0.465,0.336) (-0.450,0.354) (-0.434,0.371) (-0.418,0.388) (-0.401,0.405) (-0.383,0.422) (-0.365,0.439) (-0.347,0.456) (-0.328,0.473) (-0.309,0.489) (-0.289,0.505) (-0.269,0.520) (-0.249,0.535) (-0.229,0.549) (-0.208,0.562) (-0.186,0.575) (-0.165,0.587) (-0.144,0.597) (-0.122,0.606) (-0.100,0.614) (-0.078,0.621) (-0.056,0.626) (-0.033,0.629) (-0.011,0.631) (0.011,0.631) (0.033,0.629) (0.056,0.626) (0.078,0.621) (0.100,0.614) (0.122,0.606) (0.144,0.597) (0.165,0.587) (0.186,0.575) (0.208,0.562) (0.229,0.549) (0.249,0.535) (0.269,0.520) (0.289,0.505) (0.309,0.489) (0.328,0.473) (0.347,0.456) (0.365,0.439) (0.383,0.422) (0.401,0.405) (0.418,0.388) (0.434,0.371) (0.450,0.354) (0.465,0.336) (0.480,0.319) (0.494,0.302) (0.508,0.285) (0.521,0.268) (0.533,0.252) (0.545,0.235) (0.555,0.219) (0.566,0.202) (0.575,0.186) (0.584,0.170) (0.592,0.154) (0.599,0.138) (0.606,0.122) (0.612,0.107) (0.617,0.091) (0.621,0.076) (0.625,0.061) (0.627,0.046) (0.629,0.030) (0.631,0.015) (0.631,0.000)};
\draw[c4] plot[smooth] coordinates {(-0.631,-0.000) (-0.631,-0.015) (-0.629,-0.030) (-0.627,-0.046) (-0.625,-0.061) (-0.621,-0.076) (-0.617,-0.091) (-0.612,-0.107) (-0.606,-0.122) (-0.599,-0.138) (-0.592,-0.154) (-0.584,-0.170) (-0.575,-0.186) (-0.566,-0.202) (-0.555,-0.219) (-0.545,-0.235) (-0.533,-0.252) (-0.521,-0.268) (-0.508,-0.285) (-0.494,-0.302) (-0.480,-0.319) (-0.465,-0.336) (-0.450,-0.354) (-0.434,-0.371) (-0.418,-0.388) (-0.401,-0.405) (-0.383,-0.422) (-0.365,-0.439) (-0.347,-0.456) (-0.328,-0.473) (-0.309,-0.489) (-0.289,-0.505) (-0.269,-0.520) (-0.249,-0.535) (-0.229,-0.549) (-0.208,-0.562) (-0.186,-0.575) (-0.165,-0.587) (-0.144,-0.597) (-0.122,-0.606) (-0.100,-0.614) (-0.078,-0.621) (-0.056,-0.626) (-0.033,-0.629) (-0.011,-0.631) (0.011,-0.631) (0.033,-0.629) (0.056,-0.626) (0.078,-0.621) (0.100,-0.614) (0.122,-0.606) (0.144,-0.597) (0.165,-0.587) (0.186,-0.575) (0.208,-0.562) (0.229,-0.549) (0.249,-0.535) (0.269,-0.520) (0.289,-0.505) (0.309,-0.489) (0.328,-0.473) (0.347,-0.456) (0.365,-0.439) (0.383,-0.422) (0.401,-0.405) (0.418,-0.388) (0.434,-0.371) (0.450,-0.354) (0.465,-0.336) (0.480,-0.319) (0.494,-0.302) (0.508,-0.285) (0.521,-0.268) (0.533,-0.252) (0.545,-0.235) (0.555,-0.219) (0.566,-0.202) (0.575,-0.186) (0.584,-0.170) (0.592,-0.154) (0.599,-0.138) (0.606,-0.122) (0.612,-0.107) (0.617,-0.091) (0.621,-0.076) (0.625,-0.061) (0.627,-0.046) (0.629,-0.030) (0.631,-0.015) (0.631,-0.000)};
\draw[c5] plot[smooth] coordinates {(-0.770,0.000) (-0.769,0.014) (-0.768,0.028) (-0.766,0.043) (-0.762,0.057) (-0.758,0.072) (-0.753,0.087) (-0.747,0.102) (-0.739,0.117) (-0.731,0.132) (-0.722,0.148) (-0.713,0.164) (-0.702,0.181) (-0.690,0.198) (-0.678,0.215) (-0.665,0.233) (-0.650,0.251) (-0.635,0.269) (-0.620,0.288) (-0.603,0.308) (-0.586,0.327) (-0.568,0.347) (-0.549,0.368) (-0.530,0.388) (-0.510,0.409) (-0.489,0.431) (-0.468,0.452) (-0.446,0.474) (-0.424,0.496) (-0.401,0.518) (-0.377,0.540) (-0.353,0.562) (-0.329,0.584) (-0.304,0.606) (-0.279,0.628) (-0.253,0.648) (-0.228,0.669) (-0.201,0.688) (-0.175,0.706) (-0.149,0.722) (-0.122,0.737) (-0.095,0.749) (-0.068,0.759) (-0.041,0.766) (-0.014,0.770) (0.014,0.770) (0.041,0.766) (0.068,0.759) (0.095,0.749) (0.122,0.737) (0.149,0.722) (0.175,0.706) (0.201,0.688) (0.228,0.669) (0.253,0.648) (0.279,0.628) (0.304,0.606) (0.329,0.584) (0.353,0.562) (0.377,0.540) (0.401,0.518) (0.424,0.496) (0.446,0.474) (0.468,0.452) (0.489,0.431) (0.510,0.409) (0.530,0.388) (0.549,0.368) (0.568,0.347) (0.586,0.327) (0.603,0.308) (0.620,0.288) (0.635,0.269) (0.650,0.251) (0.665,0.233) (0.678,0.215) (0.690,0.198) (0.702,0.181) (0.713,0.164) (0.722,0.148) (0.731,0.132) (0.739,0.117) (0.747,0.102) (0.753,0.087) (0.758,0.072) (0.762,0.057) (0.766,0.043) (0.768,0.028) (0.769,0.014) (0.770,0.000)};
\draw[c5] plot[smooth] coordinates {(-0.770,-0.000) (-0.769,-0.014) (-0.768,-0.028) (-0.766,-0.043) (-0.762,-0.057) (-0.758,-0.072) (-0.753,-0.087) (-0.747,-0.102) (-0.739,-0.117) (-0.731,-0.132) (-0.722,-0.148) (-0.713,-0.164) (-0.702,-0.181) (-0.690,-0.198) (-0.678,-0.215) (-0.665,-0.233) (-0.650,-0.251) (-0.635,-0.269) (-0.620,-0.288) (-0.603,-0.308) (-0.586,-0.327) (-0.568,-0.347) (-0.549,-0.368) (-0.530,-0.388) (-0.510,-0.409) (-0.489,-0.431) (-0.468,-0.452) (-0.446,-0.474) (-0.424,-0.496) (-0.401,-0.518) (-0.377,-0.540) (-0.353,-0.562) (-0.329,-0.584) (-0.304,-0.606) (-0.279,-0.628) (-0.253,-0.648) (-0.228,-0.669) (-0.201,-0.688) (-0.175,-0.706) (-0.149,-0.722) (-0.122,-0.737) (-0.095,-0.749) (-0.068,-0.759) (-0.041,-0.766) (-0.014,-0.770) (0.014,-0.770) (0.041,-0.766) (0.068,-0.759) (0.095,-0.749) (0.122,-0.737) (0.149,-0.722) (0.175,-0.706) (0.201,-0.688) (0.228,-0.669) (0.253,-0.648) (0.279,-0.628) (0.304,-0.606) (0.329,-0.584) (0.353,-0.562) (0.377,-0.540) (0.401,-0.518) (0.424,-0.496) (0.446,-0.474) (0.468,-0.452) (0.489,-0.431) (0.510,-0.409) (0.530,-0.388) (0.549,-0.368) (0.568,-0.347) (0.586,-0.327) (0.603,-0.308) (0.620,-0.288) (0.635,-0.269) (0.650,-0.251) (0.665,-0.233) (0.678,-0.215) (0.690,-0.198) (0.702,-0.181) (0.713,-0.164) (0.722,-0.148) (0.731,-0.132) (0.739,-0.117) (0.747,-0.102) (0.753,-0.087) (0.758,-0.072) (0.762,-0.057) (0.766,-0.043) (0.768,-0.028) (0.769,-0.014) (0.770,-0.000)};
\draw[c6] plot coordinates {(-1.000,0.000) (0.000,1.000) (1.000,0.000)};
\draw[c6] plot coordinates {(-1.000,-0.000) (0.000,-1.000) (1.000,-0.000)};
\draw[c7] plot[smooth] coordinates {(0.295,1.000) (0.295,0.989) (0.296,0.977) (0.297,0.966) (0.299,0.955) (0.301,0.943) (0.304,0.932) (0.307,0.920) (0.311,0.909) (0.315,0.897) (0.320,0.885) (0.325,0.873) (0.331,0.861) (0.337,0.849) (0.344,0.837) (0.351,0.825) (0.358,0.812) (0.366,0.800) (0.375,0.787) (0.383,0.774) (0.393,0.761) (0.402,0.748) (0.413,0.734) (0.423,0.721) (0.434,0.707) (0.446,0.693) (0.457,0.679) (0.470,0.665) (0.482,0.651) (0.495,0.636) (0.508,0.622) (0.522,0.607) (0.536,0.593) (0.550,0.578) (0.565,0.563) (0.580,0.549) (0.595,0.534) (0.610,0.519) (0.626,0.504) (0.642,0.490) (0.658,0.475) (0.675,0.461) (0.692,0.447) (0.709,0.433) (0.726,0.419) (0.743,0.406) (0.761,0.393) (0.779,0.380) (0.797,0.368) (0.815,0.357) (0.833,0.346) (0.851,0.336) (0.869,0.327) (0.888,0.319) (0.906,0.312) (0.925,0.306) (0.944,0.301) (0.962,0.298) (0.981,0.296) (1.000,0.295)};
\draw[c7] plot[smooth] coordinates {(0.295,-1.000) (0.295,-0.989) (0.296,-0.977) (0.297,-0.966) (0.299,-0.955) (0.301,-0.943) (0.304,-0.932) (0.307,-0.920) (0.311,-0.909) (0.315,-0.897) (0.320,-0.885) (0.325,-0.873) (0.331,-0.861) (0.337,-0.849) (0.344,-0.837) (0.351,-0.825) (0.358,-0.812) (0.366,-0.800) (0.375,-0.787) (0.383,-0.774) (0.393,-0.761) (0.402,-0.748) (0.413,-0.734) (0.423,-0.721) (0.434,-0.707) (0.446,-0.693) (0.457,-0.679) (0.470,-0.665) (0.482,-0.651) (0.495,-0.636) (0.508,-0.622) (0.522,-0.607) (0.536,-0.593) (0.550,-0.578) (0.565,-0.563) (0.580,-0.549) (0.595,-0.534) (0.610,-0.519) (0.626,-0.504) (0.642,-0.490) (0.658,-0.475) (0.675,-0.461) (0.692,-0.447) (0.709,-0.433) (0.726,-0.419) (0.743,-0.406) (0.761,-0.393) (0.779,-0.380) (0.797,-0.368) (0.815,-0.357) (0.833,-0.346) (0.851,-0.336) (0.869,-0.327) (0.888,-0.319) (0.906,-0.312) (0.925,-0.306) (0.944,-0.301) (0.962,-0.298) (0.981,-0.296) (1.000,-0.295)};
\draw[c7] plot[smooth] coordinates {(-0.295,1.000) (-0.295,0.989) (-0.296,0.977) (-0.297,0.966) (-0.299,0.955) (-0.301,0.943) (-0.304,0.932) (-0.307,0.920) (-0.311,0.909) (-0.315,0.897) (-0.320,0.885) (-0.325,0.873) (-0.331,0.861) (-0.337,0.849) (-0.344,0.837) (-0.351,0.825) (-0.358,0.812) (-0.366,0.800) (-0.375,0.787) (-0.383,0.774) (-0.393,0.761) (-0.402,0.748) (-0.413,0.734) (-0.423,0.721) (-0.434,0.707) (-0.446,0.693) (-0.457,0.679) (-0.470,0.665) (-0.482,0.651) (-0.495,0.636) (-0.508,0.622) (-0.522,0.607) (-0.536,0.593) (-0.550,0.578) (-0.565,0.563) (-0.580,0.549) (-0.595,0.534) (-0.610,0.519) (-0.626,0.504) (-0.642,0.490) (-0.658,0.475) (-0.675,0.461) (-0.692,0.447) (-0.709,0.433) (-0.726,0.419) (-0.743,0.406) (-0.761,0.393) (-0.779,0.380) (-0.797,0.368) (-0.815,0.357) (-0.833,0.346) (-0.851,0.336) (-0.869,0.327) (-0.888,0.319) (-0.906,0.312) (-0.925,0.306) (-0.944,0.301) (-0.962,0.298) (-0.981,0.296) (-1.000,0.295)};
\draw[c7] plot[smooth] coordinates {(-0.295,-1.000) (-0.295,-0.989) (-0.296,-0.977) (-0.297,-0.966) (-0.299,-0.955) (-0.301,-0.943) (-0.304,-0.932) (-0.307,-0.920) (-0.311,-0.909) (-0.315,-0.897) (-0.320,-0.885) (-0.325,-0.873) (-0.331,-0.861) (-0.337,-0.849) (-0.344,-0.837) (-0.351,-0.825) (-0.358,-0.812) (-0.366,-0.800) (-0.375,-0.787) (-0.383,-0.774) (-0.393,-0.761) (-0.402,-0.748) (-0.413,-0.734) (-0.423,-0.721) (-0.434,-0.707) (-0.446,-0.693) (-0.457,-0.679) (-0.470,-0.665) (-0.482,-0.651) (-0.495,-0.636) (-0.508,-0.622) (-0.522,-0.607) (-0.536,-0.593) (-0.550,-0.578) (-0.565,-0.563) (-0.580,-0.549) (-0.595,-0.534) (-0.610,-0.519) (-0.626,-0.504) (-0.642,-0.490) (-0.658,-0.475) (-0.675,-0.461) (-0.692,-0.447) (-0.709,-0.433) (-0.726,-0.419) (-0.743,-0.406) (-0.761,-0.393) (-0.779,-0.380) (-0.797,-0.368) (-0.815,-0.357) (-0.833,-0.346) (-0.851,-0.336) (-0.869,-0.327) (-0.888,-0.319) (-0.906,-0.312) (-0.925,-0.306) (-0.944,-0.301) (-0.962,-0.298) (-0.981,-0.296) (-1.000,-0.295)};
\draw[c8] plot[smooth] coordinates {(0.500,1.000) (0.500,0.989) (0.501,0.979) (0.502,0.968) (0.503,0.958) (0.504,0.947) (0.506,0.936) (0.509,0.926) (0.511,0.915) (0.514,0.904) (0.518,0.894) (0.521,0.883) (0.525,0.872) (0.530,0.862) (0.534,0.851) (0.539,0.840) (0.545,0.830) (0.550,0.819) (0.556,0.808) (0.563,0.798) (0.569,0.787) (0.576,0.776) (0.583,0.766) (0.591,0.755) (0.599,0.745) (0.607,0.734) (0.615,0.724) (0.624,0.713) (0.633,0.703) (0.642,0.693) (0.651,0.683) (0.661,0.673) (0.671,0.663) (0.681,0.653) (0.691,0.643) (0.702,0.634) (0.713,0.624) (0.724,0.615) (0.735,0.606) (0.746,0.597) (0.758,0.589) (0.769,0.581) (0.781,0.573) (0.793,0.565) (0.806,0.558) (0.818,0.551) (0.830,0.544) (0.843,0.538) (0.856,0.532) (0.868,0.527) (0.881,0.522) (0.894,0.517) (0.907,0.513) (0.920,0.510) (0.934,0.507) (0.947,0.504) (0.960,0.502) (0.973,0.501) (0.987,0.500) (1.000,0.500)};
\draw[c8] plot[smooth] coordinates {(0.500,-1.000) (0.500,-0.989) (0.501,-0.979) (0.502,-0.968) (0.503,-0.958) (0.504,-0.947) (0.506,-0.936) (0.509,-0.926) (0.511,-0.915) (0.514,-0.904) (0.518,-0.894) (0.521,-0.883) (0.525,-0.872) (0.530,-0.862) (0.534,-0.851) (0.539,-0.840) (0.545,-0.830) (0.550,-0.819) (0.556,-0.808) (0.563,-0.798) (0.569,-0.787) (0.576,-0.776) (0.583,-0.766) (0.591,-0.755) (0.599,-0.745) (0.607,-0.734) (0.615,-0.724) (0.624,-0.713) (0.633,-0.703) (0.642,-0.693) (0.651,-0.683) (0.661,-0.673) (0.671,-0.663) (0.681,-0.653) (0.691,-0.643) (0.702,-0.634) (0.713,-0.624) (0.724,-0.615) (0.735,-0.606) (0.746,-0.597) (0.758,-0.589) (0.769,-0.581) (0.781,-0.573) (0.793,-0.565) (0.806,-0.558) (0.818,-0.551) (0.830,-0.544) (0.843,-0.538) (0.856,-0.532) (0.868,-0.527) (0.881,-0.522) (0.894,-0.517) (0.907,-0.513) (0.920,-0.510) (0.934,-0.507) (0.947,-0.504) (0.960,-0.502) (0.973,-0.501) (0.987,-0.500) (1.000,-0.500)};
\draw[c8] plot[smooth] coordinates {(-0.500,1.000) (-0.500,0.989) (-0.501,0.979) (-0.502,0.968) (-0.503,0.958) (-0.504,0.947) (-0.506,0.936) (-0.509,0.926) (-0.511,0.915) (-0.514,0.904) (-0.518,0.894) (-0.521,0.883) (-0.525,0.872) (-0.530,0.862) (-0.534,0.851) (-0.539,0.840) (-0.545,0.830) (-0.550,0.819) (-0.556,0.808) (-0.563,0.798) (-0.569,0.787) (-0.576,0.776) (-0.583,0.766) (-0.591,0.755) (-0.599,0.745) (-0.607,0.734) (-0.615,0.724) (-0.624,0.713) (-0.633,0.703) (-0.642,0.693) (-0.651,0.683) (-0.661,0.673) (-0.671,0.663) (-0.681,0.653) (-0.691,0.643) (-0.702,0.634) (-0.713,0.624) (-0.724,0.615) (-0.735,0.606) (-0.746,0.597) (-0.758,0.589) (-0.769,0.581) (-0.781,0.573) (-0.793,0.565) (-0.806,0.558) (-0.818,0.551) (-0.830,0.544) (-0.843,0.538) (-0.856,0.532) (-0.868,0.527) (-0.881,0.522) (-0.894,0.517) (-0.907,0.513) (-0.920,0.510) (-0.934,0.507) (-0.947,0.504) (-0.960,0.502) (-0.973,0.501) (-0.987,0.500) (-1.000,0.500)};
\draw[c8] plot[smooth] coordinates {(-0.500,-1.000) (-0.500,-0.989) (-0.501,-0.979) (-0.502,-0.968) (-0.503,-0.958) (-0.504,-0.947) (-0.506,-0.936) (-0.509,-0.926) (-0.511,-0.915) (-0.514,-0.904) (-0.518,-0.894) (-0.521,-0.883) (-0.525,-0.872) (-0.530,-0.862) (-0.534,-0.851) (-0.539,-0.840) (-0.545,-0.830) (-0.550,-0.819) (-0.556,-0.808) (-0.563,-0.798) (-0.569,-0.787) (-0.576,-0.776) (-0.583,-0.766) (-0.591,-0.755) (-0.599,-0.745) (-0.607,-0.734) (-0.615,-0.724) (-0.624,-0.713) (-0.633,-0.703) (-0.642,-0.693) (-0.651,-0.683) (-0.661,-0.673) (-0.671,-0.663) (-0.681,-0.653) (-0.691,-0.643) (-0.702,-0.634) (-0.713,-0.624) (-0.724,-0.615) (-0.735,-0.606) (-0.746,-0.597) (-0.758,-0.589) (-0.769,-0.581) (-0.781,-0.573) (-0.793,-0.565) (-0.806,-0.558) (-0.818,-0.551) (-0.830,-0.544) (-0.843,-0.538) (-0.856,-0.532) (-0.868,-0.527) (-0.881,-0.522) (-0.894,-0.517) (-0.907,-0.513) (-0.920,-0.510) (-0.934,-0.507) (-0.947,-0.504) (-0.960,-0.502) (-0.973,-0.501) (-0.987,-0.500) (-1.000,-0.500)};

  \draw[atomLine, line width=0.9pt] (-1,-1) rectangle (1,1);
  \draw[-{Stealth[length=2mm]}, fmGray] (-1.25,0) -- (1.35,0) node[right, font=\small]{$k_x$};
  \draw[-{Stealth[length=2mm]}, fmGray] (0,-1.25) -- (0,1.35) node[above, font=\small]{$k_y$};
  \draw[arr, black] (0,0) -- (0.98,0) node[pos=0.55, below, font=\small]{1};
  \draw[arr, black] (0,0) -- (0.97,0.97) node[pos=0.5, above left=-2pt, font=\small]{2};
  \fill (1,0) circle (0.7pt) node[right, font=\small]{$C$};
  \node[font=\scriptsize, below left] at (1,-1) {$\frac\pi a$};
\end{tikzpicture}
\caption{Изоэнергетические линии в первой зоне Бриллюэна квадратной
решётки (расчёт в модели сильной связи). Синие замкнутые линии около
центра~— малые энергии; красная~— линия, касающаяся границы в точках
$C$; жёлтые дуги~— высокие энергии, когда свободны лишь углы зоны}
\label{fig:contours}
\end{figure}

Дальнейшее зависит от величины разрыва энергии $\Delta E$ на границе
зоны. Максимум энергии в первой зоне находится в её углах, а минимум во
второй~— у середин сторон (точки $C$). Возможны два случая
(рис.~\ref{fig:gap-overlap}).
\begin{itemize}
\item \emph{Разрыв велик}: минимум второй зоны $E_\text{II\,min}$ лежит
\emph{выше} максимума первой $E_\text{I\,max}$. Выгодно сначала
заполнить первую зону целиком. Если электронов как раз на это хватает
(два на атом), кристалл~— \emph{диэлектрик} или \emph{полупроводник} в
зависимости от ширины щели.
\item \emph{Разрыв мал}: $E_\text{II\,min}<E_\text{I\,max}$~—
\term{перекрытие зон}. Первая зона заполняется лишь до уровня
$E_\text{II\,min}$, дальше электроны занимают состояния второй зоны
около точек $C$, а углы первой зоны остаются пустыми. Обе зоны заполнены
частично, и вещество~— \emph{металл}. Так обстоит дело у двухвалентных
металлов (Be, Mg, Zn\ldots): двух электронов на атом хватило бы, чтобы
целиком заполнить зону, но благодаря перекрытию они~— проводники.
\end{itemize}

\begin{figure}[ht]
\centering
\begin{subfigure}[b]{0.48\textwidth}\centering
\begin{tikzpicture}
\begin{axis}[width=\textwidth, height=5.4cm, xmin=0, xmax=1.48, ymin=0, ymax=2.35,
  xlabel={$k$, ед. $\pi/a$}, ylabel={$E$}, xtick={1,1.414}, xticklabels={$1$,$\sqrt2$}, ytick=\empty, axis lines=left]
  \addplot[fmRed, thick, domain=0:1, samples=60] {x^2-x^6/3};
  \addplot[fmRed, thick, domain=1:1.45, samples=40] {1.667+1.6*(x-1)^2};
  \addplot[fmBlue, thick, domain=0:1.414, samples=60] {x^2-x^6/12};
  \draw[fmBlue, dashed] (axis cs:0,1.333) -- (axis cs:1.48,1.333);
  \draw[fmRed, dashed] (axis cs:0,1.667) -- (axis cs:1.48,1.667);
  \node[font=\scriptsize, fmBlue, anchor=north west] at (axis cs:0.02,1.33) {$E_\text{I\,max}$};
  \node[font=\scriptsize, fmRed, anchor=south west] at (axis cs:0.02,1.67) {$E_\text{II\,min}$};
  \node[font=\scriptsize, fmRed] at (axis cs:0.82,0.42) {напр.~1};
  \node[font=\scriptsize, fmBlue] at (axis cs:1.2,0.95) {напр.~2};
\end{axis}
\end{tikzpicture}
\caption{щель между зонами}
\end{subfigure}\hfill
\begin{subfigure}[b]{0.48\textwidth}\centering
\begin{tikzpicture}
\begin{axis}[width=\textwidth, height=5.4cm, xmin=0, xmax=1.48, ymin=0, ymax=2.35,
  xlabel={$k$, ед. $\pi/a$}, ylabel={$E$}, xtick={1,1.414}, xticklabels={$1$,$\sqrt2$}, ytick=\empty, axis lines=left]
  \fill[fmGreen!18] (axis cs:0,1.067) rectangle (axis cs:1.48,1.333);
  \addplot[fmRed, thick, domain=0:1, samples=60] {x^2-x^6/3};
  \addplot[fmRed, thick, domain=1:1.45, samples=40] {1.067+1.6*(x-1)^2};
  \addplot[fmBlue, thick, domain=0:1.414, samples=60] {x^2-x^6/12};
  \draw[fmBlue, dashed] (axis cs:0,1.333) -- (axis cs:1.48,1.333);
  \draw[fmRed, dashed] (axis cs:0,1.067) -- (axis cs:1.48,1.067);
  \node[font=\scriptsize, fmBlue, anchor=south west] at (axis cs:0.02,1.34) {$E_\text{I\,max}$};
  \node[font=\scriptsize, fmRed, anchor=north west] at (axis cs:0.02,1.06) {$E_\text{II\,min}$};
  \node[font=\scriptsize, fmGreen!50!black] at (axis cs:0.55,1.2) {перекрытие};
  \node[font=\scriptsize, fmRed] at (axis cs:0.82,0.42) {напр.~1};
  \node[font=\scriptsize, fmBlue] at (axis cs:1.2,0.95) {напр.~2};
\end{axis}
\end{tikzpicture}
\caption{перекрытие зон}
\end{subfigure}
\caption{Энергия электрона вдоль направлений~1 (к середине стороны
зоны) и~2 (к углу зоны), схема. Кривая~1 раньше упирается в границу и
продолжается во второй зоне выше на $\Delta E$. \textit{а}~— минимум
второй зоны выше максимума первой: щель; \textit{б}~— ниже: зоны
перекрываются}
\label{fig:gap-overlap}
\end{figure}

\paragraph{Металлы, полупроводники, диэлектрики}
Почему полностью заполненная зона не проводит ток? В заполненной зоне
каждому электрону с волновым вектором $\vb k$ (и скоростью $\vb v$)
соответствует электрон с $-\vb k$ (и скоростью $-\vb v$), и суммарный
ток равен нулю. Электрическое поле не может изменить это распределение:
чтобы ускориться, электрон должен перейти в соседнее по $\vb k$
состояние, а все они заняты. Ток возможен лишь в зоне, заполненной
\emph{частично}. Отсюда классификация твёрдых тел
(рис.~\ref{fig:band-filling}).

\begin{figure}[ht]
\centering
\begin{tikzpicture}[x=1cm, y=1cm,
  bnd/.style={draw=atomLine, line width=0.5pt, fill=white},
  occ/.style={fill=fmBlue!45}]
  \foreach \s/\t/\u in {0/металл/{(Na, Cu)}, 3.4/металл/{(Mg, Zn)}, 6.8/полупроводник/{(Si, Ge)}, 10.2/диэлектрик/{(алмаз, NaCl)}}{
    \node[font=\small, align=center] at (\s+1,-0.75) {\t\\[-1pt]\footnotesize\u};}
  % натрий: зона заполнена наполовину
  \fill[occ] (0.2,0) rectangle (1.8,1.3);
  \draw[bnd, fill=none] (0.2,0) rectangle (1.8,2.6);
  \draw[fmRed, thick] (0.05,1.3) -- (1.95,1.3);
  % перекрытие: зона I (слева) и зона II (справа, сдвинута вверх)
  \fill[occ] (3.55,0) rectangle (4.75,2.0);
  \draw[bnd, fill=none] (3.55,0) rectangle (4.75,2.3);
  \fill[occ] (4.05,1.6) rectangle (5.25,2.0);
  \draw[bnd, fill=none] (4.05,1.6) rectangle (5.25,3.4);
  \node[font=\scriptsize] at (3.8,0.3) {I}; \node[font=\scriptsize] at (5.0,3.1) {II};
  \draw[fmRed, thick] (3.45,2.0) -- (5.35,2.0);
  % полупроводник
  \fill[occ] (7.0,0) rectangle (8.6,1.5);
  \draw[bnd, fill=none] (7.0,0) rectangle (8.6,1.5);
  \draw[bnd, fill=none] (7.0,2.0) rectangle (8.6,3.4);
  \draw[dimline] (8.8,1.5) -- node[right, font=\scriptsize]{$\sim1$ эВ} (8.8,2.0);
  % диэлектрик
  \fill[occ] (10.4,0) rectangle (12.0,1.0);
  \draw[bnd, fill=none] (10.4,0) rectangle (12.0,1.0);
  \draw[bnd, fill=none] (10.4,3.0) rectangle (12.0,3.4);
  \draw[dimline] (12.2,1.0) -- node[right, font=\scriptsize]{$>3$ эВ} (12.2,3.0);
  \node[font=\scriptsize, fmRed] at (1.0,1.5) {$\EF$};
\end{tikzpicture}
\caption{Заполнение зон (синее~— занятые состояния при $T=0$, красная
линия~— уровень Ферми): металл с наполовину заполненной зоной;
металл с перекрывающимися зонами; полупроводник и диэлектрик~—
заполненная зона отделена от пустой узкой или широкой щелью}
\label{fig:band-filling}
\end{figure}

\begin{itemize}
\item \term{Металлы}~— есть частично заполненная зона: либо электронов
не хватает, чтобы заполнить зону (одновалентные металлы), либо зоны
перекрываются (двухвалентные).
\item \term{Диэлектрики}~— все зоны либо полностью заполнены, либо пусты,
и щель широкая (у алмаза около 5,5~эВ).
\item \term{Полупроводники}~— то же, но щель узкая (Si~— 1,1~эВ, Ge~—
0,67~эВ). Уже при комнатной температуре тепловое движение перебрасывает
часть электронов через щель, и проводимость растёт с нагревом.
\end{itemize}

\paragraph{Поверхность Ферми}
Для металла важнее всего, как выглядит граница занятых состояний.
\begin{opredelenie}
\term{Поверхность Ферми}~— изоэнергетическая поверхность $E(\vb k)=\EF$,
отделяющая при $T=0$ занятые состояния от свободных.
\end{opredelenie}
Её форма зависит от двух факторов: структуры кристалла (формы зоны) и
числа электронов на атом.
\begin{itemize}
\item \emph{Щелочные металлы} (ОЦК). Электроны заполняют половину зоны;
радиус сферы Ферми составляет около 88\,\% расстояния до ближайших
граней (табл.~\ref{tab:bz-points}: $0{,}620/0{,}707$). Грани далеко,
поверхность Ферми почти сферическая и замкнутая; её искажения~—
двенадцать слабых выпуклостей к центрам граней~— не превышают 1,5\,\%.
\item \emph{Благородные металлы} Cu, Ag, Au (ГЦК). Зона тоже заполнена
наполовину, но сфера подходит к шестиугольным граням $\{111\}$ гораздо
ближе~— на $0{,}782/0{,}866\approx90$\,\%. К тому же потенциал ионных
остовов сильнее, чем у щелочных металлов. Опыт показывает, что
поверхность Ферми вытягивается и \emph{касается} шестиугольных граней
«шейками» (рис.~\ref{fig:necks}). Поверхность становится
\emph{открытой}: в повторяющихся зонах она тянется через всё
$\vb k$-пространство. Диаметр шеек сильно зависит от потенциала решётки.
\item \emph{Многовалентные и переходные металлы}~— поверхности Ферми
сложные. У алюминия (ГЦК, три электрона на атом) можно было бы
ожидать заполненной первой зоны и половины второй. Расчёт показывает,
что электроны занимают состояния \emph{четырёх} зон: первая заполнена
целиком, а куски поверхности Ферми лежат во второй, третьей и
четвёртой зонах.
\end{itemize}

\begin{figure}[ht]
\centering
\begin{tikzpicture}[scale=0.9]
  \foreach \cx in {0,3.0,6.0}{
    \draw[atomLine, line width=0.6pt] (\cx,0) ++(30:1.732) \foreach \a in {90,150,...,390}{-- ($(\cx,0)+(\a:1.732)$)};
  }
  \foreach \cx in {1.5,4.5}{
    \fill[fmBlue!35] (\cx-0.35,-0.17) rectangle (\cx+0.35,0.17);}
  \foreach \cx in {0,3.0,6.0}{
    \fill[fmBlue!35] (\cx,0) circle (1.38);
    \draw[fmBlue!70!black, line width=0.6pt] (\cx,0) circle (1.38);}
  \foreach \cx in {1.5,4.5}{
    \fill[fmBlue!35] (\cx-0.2,-0.17) rectangle (\cx+0.2,0.17);
    \draw[fmBlue!70!black, line width=0.6pt] (\cx-0.16,0.17) -- (\cx+0.16,0.17) (\cx-0.16,-0.17) -- (\cx+0.16,-0.17);}
  \draw[fmRed, line width=0.5pt] (4.5,0.2) -- (5.3,1.95) node[above, font=\small]{шейка};
  \node[font=\small] at (0,-2.05) {зона 1};
  \node[font=\small] at (3,-2.05) {соседняя зона};
  \node[font=\small] at (6,-2.05) {\ldots};
\end{tikzpicture}
\caption{Открытая поверхность Ферми благородных металлов (схема сечения
в повторяющихся зонах): почти сферические поверхности соседних зон
соединены шейками через шестиугольные грани}
\label{fig:necks}
\end{figure}

\begin{smysl}
Форма (топология) поверхности Ферми сильно влияет на свойства металла.
А она может меняться: тепловое возбуждение, дефекты, примеси и
легирование меняют число электронов на атом и строение зон. Это одна из
причин, по которым свойства сплавов так чувствительны к составу.
\end{smysl}

\paragraph{Плотность состояний в зонной теории}
В теории свободных электронов $g(E)\propto\sqrt E$ (рис.~\ref{fig:fermi-dirac}).
В зонной теории кривая меняется (рис.~\ref{fig:dos-band},\,\textit{а}).
\begin{itemize}
\item Участок $OA$: заполняются состояния с низкой энергией вблизи
центра зоны, электроны почти свободны, кривая совпадает с параболой.
\item Участок $AB$: заполняются состояния вблизи границ зоны, где
энергия растёт с $k$ очень медленно, поэтому на единичный интервал
энергии приходится больше состояний: $g(E)$ идёт выше параболы.
\item Точка $B$: поверхность Ферми достигает границы зоны и меняет
связность; производная $\d g/\d E$ терпит разрыв (излом кривой).
\item Участок $BCD$: остаются незаполненными лишь углы зоны,
состояний становится всё меньше, и $g(E)$ падает до нуля в точке $D$,
где зона заполнена полностью.
\end{itemize}
Если зоны перекрываются, общая плотность состояний~— сумма кривых
отдельных зон (рис.~\ref{fig:dos-band},\,\textit{б}): часть электронов с
верхних уровней первой зоны переходит на более низкие свободные уровни
второй, и обе зоны заполнены частично. Если не перекрываются~— у каждой
зоны своя кривая, а между ними щель (\textit{в}). Точный расчёт $g(E)$
для реальных металлов~— сложная задача.

\begin{figure}[ht]
\centering
\begin{subfigure}[b]{0.4\textwidth}\centering
\begin{tikzpicture}
\begin{axis}[width=\textwidth, height=5.2cm, xmin=-6.3, xmax=6.4, ymin=0, ymax=0.17,
  xlabel={$E$}, ylabel={$g(E)$}, xtick=\empty, ytick=\empty, axis lines=left]
  \addplot[fmGray, dashed] table {figures/dos-free.dat};
  \addplot[fmBlue, thick] table {figures/dos-sc.dat};
  \node[font=\scriptsize, below] at (axis cs:-6,0) {$O$};
  \node[font=\scriptsize, above left] at (axis cs:-4.6,0.07) {$A$};
  \node[font=\scriptsize, above] at (axis cs:-2,0.148) {$B$};
  \node[font=\scriptsize, above] at (axis cs:2,0.148) {$C$};
  \node[font=\scriptsize, below] at (axis cs:6,0) {$D$};
\end{axis}
\end{tikzpicture}
\caption{}
\end{subfigure}\hfill
\begin{subfigure}[b]{0.29\textwidth}\centering
\begin{tikzpicture}
\begin{axis}[width=1.3\textwidth, height=5.2cm, xmin=0, xmax=2.1, ymin=0, ymax=3.2,
  xlabel={$E$}, xtick=\empty, ytick=\empty, axis lines=left]
  \addplot[fmGray!70, dashed] table[x=E, y=g1] {figures/dos-overlap.dat};
  \addplot[fmGray!70, dashed] table[x=E, y=g2] {figures/dos-overlap.dat};
  \addplot[fmBlue, thick] table[x=E, y=s] {figures/dos-overlap.dat};
\end{axis}
\end{tikzpicture}
\caption{}
\end{subfigure}\hfill
\begin{subfigure}[b]{0.29\textwidth}\centering
\begin{tikzpicture}
\begin{axis}[width=1.3\textwidth, height=5.2cm, xmin=0, xmax=2.3, ymin=0, ymax=3.2,
  xlabel={$E$}, xtick=\empty, ytick=\empty, axis lines=left]
  \addplot[fmBlue, thick] table[x=E, y=s] {figures/dos-gap.dat};
  \draw[dimline] (axis cs:1.0,0.25) -- node[above, font=\scriptsize]{щель} (axis cs:1.25,0.25);
\end{axis}
\end{tikzpicture}
\caption{}
\end{subfigure}
\caption{Плотность состояний в зонной теории (расчёт в модели сильной
связи для простой кубической решётки): \textit{а}~— одна зона, пунктир~—
парабола свободных электронов; \textit{б}~— две перекрывающиеся зоны
(пунктир~— вклады зон, сплошная~— сумма); \textit{в}~— зоны разделены
щелью}
\label{fig:dos-band}
\end{figure}

\begin{itogi}
\begin{itemize}
\item В периодическом поле волновая функция электрона~— функция Блоха
$u_{\vb k}(\vb r)e^{i\vb k\vb r}$; электрон принадлежит всему кристаллу,
идеальная решётка его не рассеивает.
\item Разрешённые энергии образуют зоны, разделённые запрещёнными
щелями. Щели возникают при брэгговском отражении электронных волн:
в одномерном случае при $k=n\pi/a$, в общем~— на плоскостях, делящих
пополам векторы обратной решётки.
\item Зоны Бриллюэна: куб (простая кубическая), усечённый октаэдр (ГЦК),
ромбический додекаэдр (ОЦК). Каждая вмещает два электрона на атом.
\item Одновалентные металлы заполняют половину зоны; у двухвалентных
металлов зоны перекрываются. Частично заполненная зона~— металл;
заполненная зона и щель~— полупроводник или диэлектрик.
\item Поверхность Ферми щелочных металлов почти сферическая, у
благородных металлов~— открытая (шейки), у многовалентных~— сложная.
\end{itemize}
\end{itogi}

\begin{proverka}
\item Что общего у функции Блоха с плоской волной и чем она отличается?
\item Почему при $k=\pi/a$ возникает разрыв энергии? Нарисуйте
распределение электронной плотности в двух стоячих волнах.
\item Найдите расстояние от центра до граней первой зоны Бриллюэна для
ОЦК-решётки с $a=0{,}43$~нм (натрий).
\item Почему двухвалентный магний~— металл, хотя двух электронов на атом
хватило бы, чтобы целиком заполнить зону?
\item Почему заполненная зона не даёт тока?
\end{proverka}
